%% Clear sans, courier

% Klasa:
\documentclass[12pt]{mwart}

% Kodowanie:
\usepackage[utf8]{inputenc}

% odstep w listach
\usepackage[titles]{tocloft}
\addtolength{\cftfignumwidth}{10pt}

% Czcionki:
\usepackage[sfdefault]{ClearSans} %% option 'sfdefault' activates Clear Sans as the default text font
\usepackage[T1]{fontenc}
\usepackage{courier}
\usepackage[T1]{polski}

% Znak euro
\usepackage[official]{eurosym}

% Geometria i style
%\usepackage{showframe} % debug
\usepackage{geometry}
\geometry{a4paper,twoside=true,top=20mm,bottom=30mm,inner=35mm,outer=20mm}
\pagestyle{plain}
\usepackage{float}
\usepackage{afterpage}

% Grafika
\usepackage{graphicx}
\usepackage{epstopdf}

% matma
\usepackage{amssymb,amsmath}

% theoremy
\usepackage{amsthm}

% Stopnie
\usepackage{gensymb}

% Opisy
\usepackage{caption}
\captionsetup[table]{singlelinecheck=off}

% Tabele
\usepackage{float}
\floatstyle{plaintop}
\restylefloat{table}
\usepackage{array}
\usepackage{makecell}
\usepackage{color}

% Footnoty
\usepackage{footnote}
\usepackage{perpage} %the perpage package
\MakePerPage{footnote} %the perpage package command

% kody
\usepackage{lipsum}
\usepackage{listings}

% kolory
\definecolor{mygreen}{rgb}{0,0.6,0}
\definecolor{mygray}{rgb}{0.5,0.5,0.5}
\definecolor{mymauve}{rgb}{0.58,0,0.82}	

\lstset{ %
  backgroundcolor=\color{white},   			% choose the background color
  basicstyle=\ttfamily\footnotesize,        % size of fonts used for the code
  breaklines=true,                 			% automatic line breaking only at whitespace
  captionpos=t,                    			% sets the caption-position to top
  commentstyle=\color{mygreen},    			% comment style
  escapeinside={\%*}{*)},          			% if you want to add LaTeX within your code
  keywordstyle=\color{blue},       			% keyword style
  stringstyle=\color{mymauve},     			% string literal style
  numbers=left,					   			% line numbers on
  frame=/lines,								% frame box config
}

% Numerowanie obiektow wedlug rozdzialow
\usepackage{chngcntr}
\AtBeginDocument{\counterwithin{lstlisting}{section}}
\counterwithin{figure}{section}
\counterwithin{table}{section}
\numberwithin{equation}{section}


% Linki
\usepackage{url}
\usepackage{fancyref}

% Listy
\usepackage[ampersand]{easylist}
\ListProperties(Hide=100, Hang=true, Progressive=3ex, Style*=-- ,
Style2*=$\bullet$ ,Style3*=$\circ$ ,Style4*=\tiny$\blacksquare$ )

% Dane:
\title{Uniwersalny inteligentny sterownik silników do zastosowań przemysłowych}
\author{Łukasz Przeniosło}

% Zmienne srodowiskowe:
\newcommand{\paperTitlePl}{UNIWERSALNY INTELIGENTNY STEROWNIK SILNIKÓW DO ZASTOSOWAŃ PRZEMYSŁOWYCH}
\newcommand{\paperTitleEn}{UNIVERSAL INTELIGENT MOTOR DRIVER FOR INDUSTRIAL APPLICATIONS}
\newcommand{\topicChooseDate}{01.08.2015 r.}
\newcommand{\topicAcceptationDate}{01.06.2016 r.}
\newcommand{\tocLessLeft}[1]{\textsc{\textbf{\large \flushleft{#1}}} \vspace{3mm} \\ \indent }
\newcommand{\tocLessLeftNorm}[1]{\textsc{\textbf{\normalsize \flushleft{#1}}} \vspace{0mm} \\ \indent }
\newcommand{\tocLessCenter}[1]{\textsc{\textbf{\large \centerline{#1}}} \vspace{3mm} \\ \indent }
\newcommand{\insertImg}[3]{%
\begin{figure}[h]
    \centering
    \includegraphics[height=50mm]{#1}
    \caption{#2} \label{#1}
    \raggedright Źródło: Na podstawie \cite{#3}
\end{figure}
}
\newcommand{\insertImgSetSize}[4]{%
\begin{figure}[h]
    \centering
    \includegraphics[height=#2mm]{#1}
    \caption{#3}\label{#1}
    \raggedright Źródło: Na podstawie \cite{#4}
\end{figure}
}
\newcommand{\insertTab}[5]{%
\begin{savenotes}
    \begin{table}[h]
        \raggedright Źródło: Na podstawie \cite{#4}\\
        \vspace{2mm}
        \centering
        \begin{tabular}{#1} 
        #2
        \end{tabular}
        \caption{#3}  
        \label{#5}
        \floatstyle{plaintop}
    \end{table}
\end{savenotes}
}
\newcommand{\insertCode}[4]{%
\lstinputlisting[language=#2,caption={#3},label=#1]{#1}
Źródło: Na podstawie \cite{#4} \\
}

% Zmienic jesli jednak Engineering -.- ...
\newcommand{\firma}{{Mechatronic Systems}}

% twierdzenia
\newtheorem{defn}{Definicja}[section]

% Indeksowanie w pdfie wyjsciowym
\usepackage{hyperref}

% Naprawienie referencji do figur i tabel
\renewcommand{\theHfigure}{figure.section.thesection.\thefigure}

% START:
\begin{document}
\thispagestyle{empty} % usun nr strony ze strony tytulowej

\begin{figure}[h]
	\centering
	\includegraphics[height=18mm]{grafiki/logo_zut.pdf}
	\hspace{5mm}
	\includegraphics[height=18mm]{grafiki/logo_zut_we.pdf} \\
	\vspace{20mm}


	{\bf \large Łukasz Przeniosło} \\ 
	\vspace{5mm}

	nr albumu: 21789 \\
	kierunek studiów: Elektrotechnika \\
	specjalność: urządzenia i instalacje elektryczne \\
	forma studiów: studia niestacjonarne \\
	\vspace{20mm}

	\textsc{\textbf{\large \paperTitlePl}} \\
	\vspace{5mm}

	{\bf \large \paperTitleEn} \\
	\vspace{45mm}
	
	
	Praca dyplomowa magisterska \\
	napisana pod kierunkiem: \\
	\vspace{2mm}
	
	{\bf \large dra inż. Michała Bonisławskiego} \\
	\vspace{2mm}
	
	Katedra Elektroenergetyki i Napędów Elektrycznych
	\vspace{5mm}
	
\end{figure}

\begin{flushleft}
	{\footnotesize Data wydania tematu pracy: \topicChooseDate} \\
	{\footnotesize Data złożenia tematu pracy: \topicAcceptationDate} 
\end{flushleft}

\vspace{\fill \centerline{Szczecin 2016}}
\clearpage

\afterpage{\null \thispagestyle{empty} \newpage}
\clearpage

% Oswiadczenie autora
\thispagestyle{empty}

\tocLessCenter{OŚWIADCZENIE AUTORA PRACY DYPLOMOWEJ} 

Oświadczam, że praca dyplomowa inżynierska/magisterska pt. \\
 „\paperTitlePl” \\
napisana pod kierunkiem: dr. inż. Michała Bonisławskiego \\
jest w całości moim samodzielnym autorskim opracowaniem sporządzonym przy wykorzystaniu wykazanej w pracy literatury przedmiotu i materiałów źródłowych.
Złożona w dziekanacie Wydziału Elektrycznego treść mojej pracy dyplomowej w formie elektronicznej jest zgodna z treścią w formie pisemnej. \\

Oświadczam ponadto, że złożona w dziekanacie praca dyplomowa ani jej fragmenty nie
były wcześniej przedmiotem procedur procesu dyplomowania związanych z uzyskaniem
tytułu zawodowego w uczelniach wyższych.
\vspace{30mm}


\begin{flushright}
	\begin{tabular}{c}
  		............................................\\
  		podpis dyplomanta
	\end{tabular}
\end{flushright}

\vspace{5mm}

\begin{flushleft}
	Szczecin, dn. \topicAcceptationDate{}
\end{flushleft}


\clearpage


% Streszczenie
\thispagestyle{empty}

\tocLessLeft{Streszczenie pracy}
Praca przedstawia projekt systemu sterowniczego do różnorodnych silników elektrycznych, na który składają się między innymi: Zbudowane w oparciu o 32 bitowy mikrokontroler urządzenie do którego podłączany jest silnik i peryferia (sterownik) oraz aplikacja PC-owa służąca do nastaw niezbędnych do poprawnego działania układu parametrów brzegowych oraz badania pracy zasilanego przez sterownik silnika. Rozwiązanie ma zostać zastosowane do sterowania elementami ruchomymi w maszynach typu {\it Pick and Place}, drukarkach szablonowych, drukarkach 3D oraz we wszelkiego rodzaju konwojerach. W części opisowej dokumentu, została zaprezentowana niezbędna teoria nt. silników elektrycznych oraz sterowania. Ponadto, zostały porównane komercyjne rozwiązania sterowników różnej klasy, które są dostępne na rynku.

\tocLessLeft{Słowa kluczowe}
Sterowanie, silnik, elektronika mocy, oprogramowanie

\tocLessLeft{Abstract}
The paper presents a project of a control system, that can be used with a various types of electric motors. Among others the main components of the system are: built upon a 32 bit microcontroller device, to which the motors and peripherals are connected (the driver) and a PC application used for defining critical operating parameters and for gathering connected motor work diagnostics . This application is to be used with Pick and Place machines, stencil printers, 3D printers and all sort of conveyors. In the generic part of the text, the essential theory of electric motors and control has been described. Also, commercially available on the market  motor drivers of various class has been compared.

\tocLessLeft{Key words}
Controll, motor, power electronics, software


\clearpage


% Spis tresci
\tableofcontents
\clearpage



% Wstep
\section*{Wprowadzenie}

Chcąc zaprojektować system napędowy dla określonej aplikacji należy wziąć pod uwagę bardzo wiele czynników. Najważniejsze z nich to: koszt i dostępność rozwiązań gotowych na rynku. Ceny rozwiązań przemysłowych bardzo rzadko mają odzwierciedlenie w zastosowanym w ich strukturze {\it hardwarze}, jest to raczej cena za {\it know-how} producenta, wiedzę ekspercką, wsparcie techniczne oraz dziesiątki roboczogodzin programistów i elektroników tworzących produkt. W przypadku posiadania dużych środków takie rozwiązanie wydaje się być najbardziej odpowiednie. Ma ono jednak jedną znaczącą wadę (pomijając cenę oczywiście). Pomimo tego że producenci sterowników do silników prześcigają się w ilości posiadanych cech, czy technologii, to korzystając z takiego urządzenia nijako jesteśmy zmuszeni z korzystania z zamkniętego, ograniczonego systemu. System taki może lepiej lub gorzej wpasowywać się w ramy naszej aplikacji. Rzadko kiedy wpasowuje się idealnie, a najczęściej brakuje mu jakichś właściwości, lub jest przewymiarowany. \

Drugą opcją jest zakup taniego zintegrowanego chipu, który po raz kolejny w zależności od wyłożonych środków, (w o wiele mniejszej skali) zawiera w sobie mniej lub więcej logiki i czasami końcówki mocy. Taki układ można także zbudować z elementów dyskretnych, co nie obniża znacząco kosztów, lecz daje możliwość większej kontroli nad parametrami układu. Wady takiego rozwiązania są takie, że zintegrowane sterowniki są projektowane pod zastosowanie z konkretnym rodzajem silnika (np tylko do silników krokowych) i służą do sterowania silnikami małej mocy jeśli mają zintegrowane końcówki. W przypadku stosowania układów tego typu oraz dyskretnych, projektant jest zmuszony zaprojektować obwód drukowany oraz napisać bardzo dużą ilość kodu zawierającego wiele skomplikowanych procedur, które muszą działać w czasie rzeczywistym. \

W przypadku, kiedy ilość produkowanych maszyn współpracującymi ze sterownikami do silników jest duża, posiadając odpowiednią wiedzę i umiejętności, bardziej opłaca się zaprojektować rozwiązanie własne, "od zera". Ponadto, proces tworzenia tego typu systemu dostarcza ogromną ilość wiedzy praktycznej.
  
\subsection*{Cel pracy}

Zbudować sterownik do sterowania różnego typu silników (krokowych, {\it BLDC}, liniowych i {\it DC}), z możliwością dostosowywania parametrów, używając aplikacji PC. Aplikacja ma pozwalać na strojenie działających w sterowniku regulatorów, wybierać tryby pracy i badać działanie systemu w czasie rzeczywistym. Urządzenie ma zostać użyte w maszynach przemysłowych służących do produkcji urządzeń elektronicznych.

\subsection*{Zakres pracy}

Przegląd i analiza rozwiązań technologicznych, układów elektronicznych i oprogramowania stosowanych w sterownikach silników oferowanych przez wybranych producentów. Projekt uniwersalnego, kompaktowego sterownika do silników {\it DC} (prądu stałego), {\it BLDC (Brushless DC Motor)}, krokowych oraz liniowych {\it VCM (Voice Coil Motor}, inaczej {\it VCA - Voice Coil Actuator}) wraz z opisem sposobów regulacji i algorytmów zawierających elementy {\it AI}, które zostaną zaimplementowane w nowoczesnym 32 bitowym mikrokontrolerze z rodziny {\it ARM}. Opracowanie aplikacji na PC umożliwiającej dostrajanie parametrów dla konkretnego silnika oraz analizowanie jego pracy. Wykonanie, uruchomienie i badania laboratoryjne prototypu urządzenia. Prezentacja przykładowego zastosowania sterownika działającego w rzeczywistym układzie.

\clearpage


%Start czesci teoretycznej
\section{Elektromagnetyzm}

Aby zrozumieć zasadę działania większości silników elektrycznych, należy najpierw zaznajomić się z podstawowymi jednostkami występującymi w elektromagnetyzmie, dzięki którym silniki działają. Rozdział został opracowany na podstawie źródła \cite{jaszczuk}

\subsection{Natężenie pola magnetycznego}

\insertImg{grafiki/right_hand_rule.pdf}
	   {Reguła prawej dłoni - sposób wyznaczania zwrotu linii sił pola magnetycznego wokół przewodnika w którym jest prąd elektryczny}
	   {wikiRegPrawDlon}








% Opis silnikow (podrozdzialy beda includowane w rozdzialach)
\section{Podstawowe informacje o silnikach}

Rozdział ten przedstawia podział maszyn elektrycznych ze względu na ich różne parametry, oraz opis tych silników, które są (i mogą być) wspierane przez skonstruowany sterownik. 
Rozdział został opracowany głównie na podstawie źródeł \cite{przepiorkowski} i \cite{jones}.

\subsection{Podział silników elektrycznych}

Istnieje bardzo wiele rodzajów silników, a ponadto cały czas powstają nowe konstrukcje, które trudno jest sklasyfikować przy pomocy klasycznych definicji. Podział na silniki AC i DC nie wystarczy-- w przypadku zastosowania elektronicznego modułu sterującego rodzaj zasilania silnika nie ma tak dużego znaczenia. Nie można także jednoznacznie stwierdzić, czy maszyna z wirującymi magnesami to silnik prądu stałego, czy przemiennego. Biorąc to pod uwagę, silniki elektryczne można podzielić na trzy podstawowe grupy.

\tocLessLeftNorm{Silniki prądu stałego -- DC} 

Zaliczane są tu silniki komutatorowe:
\begin{easylist}
	& z magnesem trwałym (PMDC),
	& szeregowe,
	& bocznikowe,
	& szeregowo--bocznikowe.
\end{easylist} 

\tocLessLeftNorm{Silniki prądu przemiennego -- AC}

Do tej grupy  zaliczane są silniki zasilane bezpośrednio z sieci prądu przemiennego. W tym przypadku podział jest bardziej skomplikowany, ze względu na inne kryteria:
\begin{easylist}
	& \underline{Silniki asynchroniczne} -- charakteryzuje je tzw. ,,poślizg''. Rotor takiego silnika wiruje z niższą prędkością od prędkości wirowania pola magnetycznego. Najpopularniejszymi silnikami z tej grupy są silniki klatkowe.
	& \underline{Silniki synchroniczne} -- rotor wiruje z prędkością równą prędkości wirowania pola magnetycznego. Mogą być jednofazowe lub trójfazowe. Ze względu na konstrukcję, silniki te dzielą się na reluktancyjne lub z wirującym magnesem.
	& \underline{Silniki komutatorowe uniwersalne} -- komutatorowe silniki szeregowe przystosowane konstrukcyjnie do zasilania z jednej fazy prądu przemiennego.
\end{easylist} 

\tocLessLeftNorm{Silniki z komutacją elektroniczną}

Do tej grupy zalicza się silniki, których działanie nie było by możliwe bez elektronicznego układu sterującego. Najpopularniejsze z nich to:
\begin{easylist}
	& \underline{Silniki krokowe} -- ze względu na konstrukcję wyróżnia się silniki:
		&& o zmiennej reluktancji,
		&& z magnesem trwałym,
		&& hybrydowe.
\end{easylist} 
Ze względu na sposób sterowania uzwojeniami:
\begin{easylist}
		&& silniki unipolarne,
		&& silniki bipolarne.
\end{easylist} 
	Daje to sześć różnych kombinacji konstrukcyjnych silników krokowych.
	
\insertImg{grafiki/stepper_nema17.png}
	{Przykładowy silnik krokowy hybrydowy ST4118 w standardzie Nema17, produkowany przez firmę Nanotec}
	{nanotec}	
	
Nema ({\em National Electrical Manufacturers Association}) to stowarzyszenie producentów urządzeń elektrycznych, którego siedziba mieści się w USA. Standard Nema określa między innymi wymiary silników krokowych. Numer 17 oznacza szerokość i wysokość silnika 1.7 x 1.7 cala (43.2 x 43.2 mm).

\insertImg{grafiki/bldc_db42.png}
	{Przykładowy silnik BLDC (model DB42), produkowany przez firmę Nanotec}
	{nanotec}

\begin{easylist}
	& \underline{Silniki bezszczotkowe z wirującym magnesem (BLDC)} -- w literaturze angielskiej silniki te są opisywane jako prądu stałego ({\em BrushLess DC}). W przypadku silników dwubiegunowych z czujnikami Halla (np. tego typu, które są używane w małych wentylatorach do chłodzenia radiatorów) taki opis może być uzasadniony. Nie można jednak jednoznacznie zaliczyć do grupy DC silników wielobiegunowych trójfazowych z wirującym magnesem \cite{przepiorkowski}. Mylne może być także to, że w przypadku silników z wbudowanym sterownikiem nie różnią się one z punktu użytkownika od silnika komutatorowego z magnesem trwałym. Maszyny BLDC charakteryzują się bardzo wysoką prędkością obrotową, w szczególności w porównaniu do silników krokowych.	
	
	& \underline{Silniki liniowe} -- wynikiem działania takiej maszyny jest przemieszczenie liniowe, a nie ruch obrotowy. Konstrukcja takiego silnika w dużym uproszczeniu polega na ,,rozwinięciu'' wirnika i stojana na płaszczyźnie. Nazwa silnik liniowy lub {\em Linear Actuator} bardzo często dotyczy zwykłego silnika krokowego z wbudowaną przekładnią śrubową zamieniającą ruch obrotowy na posuwisto--zwrotny, co bardzo często wprowadza niejasności i wątpliwości podczas projektowania systemów napędowych.
	
	\insertImg{grafiki/linear_actuator_l35.png}
		{Przykładowy silnik krokowy hybrydowy L35 z przekładnią śrubową, produkowany przez firmę Nanotec}
		{nanotec}
	
	& \underline{Silniki liniowe VCM} -- zasada działania silnika VCM ({\em Voice Coil Motor}) lub VCA ({\em Voice Coil Actuator}) jest bardzo zbliżona do zasady działania cewki głośnikowej w szczelinie magnesu. Przesunięcie liniowe tych silników nie przekracza najczęściej 50 mm.		
	
\end{easylist} 

\insertImg{grafiki/vcm_moticont.png}
	{Przykładowy silnik liniowy VCM (model GVCM-051-051-01), produkowany przez firmę Moticont}
	{moticont}
	
\subsection{Parametry silników}

Podstawowe parametry silnika są czasami podawane na jego tablicy znamionowej. Pozostałe należy odszukać w jego karcie katalogowej lub oszacować według właściwości poszczególnych rodzajów silników.
\\

\begin{easylist}

	& Moc znamionowa: podawana jest moc elektryczna w watach pobierana przez silnik przy pracy z normalną prędkością obrotową. Moc mechaniczna jest mniejsza i zależy od sprawności silnika (typowo od 40 do 80 \%).
	\\
	& Napięcie zasilania: znamionowa wartość napięcia zasilającego (stałego lub zmiennego), przy której określane są pozostałe parametry. Najczęściej silniki mogą być zasilane napięciem niższym. Nie należy stosować napięć dużo wyższych od znamionowego (maks. +15 \%).
	\\
	& Moment obrotowy: parametr ten jest podawany w niutonometrach $ (Nm) $. Typowe wartości dla małych silników: od 0.01 do 50 Nm. Dla większości rodzajów silników występuje zależność momentu obrotowego i natężenia prądu pobieranego przez maszynę.
	\\
	& Moment rozruchowy: ważny parametr informujący o tym czy silnik jest w stanie wystartować pod obciążeniem. Podawany w Nm lub w procentach momentu obrotowego. W zależności od rodzaju silnika moment może być mały (do 150 \%), średni (150 - 250 \%) lub duży (> 250 \%).
	\\
	& Obroty znamionowe: wyrażane w obrotach na minutę ($ RPM $ -- {\em Revolutions per minute}. Zawsze podawane są przy znamionowym obciążeniu i napięciu zasilania lub podana jest charakterystyka obrotów w funkcji obciążenia. Obroty silników zawierają się w granicach od 100 do 100000 RPM, ale najczęściej spotykane wartości to od 1000 do 4000 RPM. Dla maszyn prądu przemiennego podawane są obroty przy określonej częstotliwości napięcia sieciowego (50 lub 60 Hz). Dla silników krokowych obroty znamionowe nie są podawane. Zamiast tego występują obroty maksymalne lub dopuszczalna częstotliwość impulsów.
	\\
	& Prąd znamionowy i rozruchowy: pobór prądu w normalnych warunkach pracy i w czasie rozpędzania silnika. Ten drugi może być nawet 2 - 8 krotnie większy od znamionowego (największe prądy rozruchowe mają silniki indukcyjne). Prąd rozruchowy nie występuje w silnikach sterowanych impulsowo.
	\\
	& Kierunek obrotów: parametr dotyczący przede wszystkim silników komutatorowych szeregowych i niektórych silników indukcyjnych jednofazowych. Dla tych silników kierunek obrotów jest wymuszony przez konstrukcję silnika. \newline 
	W pozostałych maszynach kierunek zależy od polaryzacji przyłożonego do ich zacisków napięcia lub w przypadku maszyn ze sterowaniem impulsowym-- od kolejności podawanie impulsów na poszczególne uzwojenia.
	\\
	& Konstrukcja mechaniczna: wymiary, ciężar, średnica wału, mocowania, rodzaj łożysk wirnika itp.
	\\
	& Konstrukcja obudowy, chłodzenie: spotykane są silniki o konstrukcji otwartej, w której uzwojenia są dostępne z zewnątrz lub zamkniętej w której silnik jest całkowicie obudowany. W silnikach otwartych często wymagane jest chłodzenie strumieniem powietrza-- np. silniki odkurzaczy, czy wiertarek.  

\end{easylist} 

\subsection{Silnik liniowy (VCM)}

Pierwszym z silników obsługiwanych przez projektowany sterownik jest silnik liniowy VCM. Jego konstrukcja jest bardzo prosta, aczkolwiek sam proces technologiczny wymagany do jego wytworzenia już nie. Z tego właśnie powodu VCM nawet bardzo małej mocy są drogie (Cena od \$100 za pojedyncze sztuki). Ciężko jest tak naprawdę stwierdzić czy jest to silnik czy siłownik, bo zakres jego ruchu nie przekracza zwykle kilku centymetrów.

\insertImg{grafiki/vcm_budowa.png}
	{Budowa cylindrycznego silnika VCM}
	{przepiorkowski}

Nazwa silnika pochodzi od jego zasady działania- identycznie jak w zwykłym głośniku jest to cewka poruszająca się w szczelinie magnesu. Zgodnie z regułą Lorentza:

\begin{defn}
	Jeżeli przez cewkę znajdującą się w polu magnetycznym przepływa prąd, to na cewkę działa siła proporcjonalna do natężenia prądu, a napięcie samoindukcji jest proporcjonalne do szybkości poruszania się cewki. \cite{przepiorkowski}.
\end{defn}

Podstawowe części składowe silnika VCM to:
\begin{easylist}
	& dwa magnesy trwałe ułożone tak, aby były skierowane tym samym biegunem w stronę cewki,
	& cewka poruszająca się w szczelinie między magnesami,
	& obudowa z miękkiego żelaza zamykająca obwód magnetyczny.
	\\
\end{easylist}

Zmieniając natężenie i kierunek przepływającego prądu, możliwe jest bardzo precyzyjne sterowanie położeniem cewki. Z tego powodu układ elektroniczny sterujący silnikiem nie wymaga wielu komponentów, a algorytm sterowania nie jest skomplikowany to VCM ma jednak jeszcze jedną zaletę, która sprawia że jest on wybierany do niektórych aplikacji zamiast np. tańszego silnika krokowego z przekładnią śrubową- pozwala wykonywać bardzo dynamiczne ruchy i posiada wysoką zwrotność. ,,Liniowy'' silnik krokowy o podobnej mocy nawet z zastosowaniem śruby o dużym skoku, nie będzie w stanie osiągnąć podobnych prędkości przy zachowaniu korzystnego momentu. Brak konieczności zamiany ruchu obrotowego na posuwisto-zwrotny eliminuje błędy pozycji (które mogą wynikać np. z luzów w przekładni). Silniki VCM stosowane są wszędzie tam, gdzie potrzebne jest dokładne i szybkie ustawienie pozycji, np. w precyzyjnych urządzeniach mechanicznych i optycznych. \\

Aby efektywnie i bezpiecznie korzystać z silników VCM w danej aplikacji, należy mieć na uwadze kilka ważnych cech tych urządzeń:

\begin{easylist}
	& Korpus cewki nie jest w żaden sposób przymocowany do obudowy silnika. Oznacza to że to projektant musi zadbać o to, aby korpus wsuwał i wysuwał się z cylindra równolegle, stosując odpowiednie mocowanie mechaniczne obu części. Nie zapewniwszy tego, uzwojenie silnika jest narażone na starcie się lakieru izolacyjnego, poprzez ocieranie się o obudowę. 
	
	& W zależności od aktualnego położenia (wysunięcia) korpusu, zależność między siłą a przepływającym przez uzwojenie prądem jest zmienna.
	
	\insertImgSetSize{grafiki/vcm_force_current_plot.eps}
		{80}
		{Przykładowy przebieg siły $ [\frac{N}{A}] $ pchającej/ ciągnącej w funkcji wysunięcia korpusu dla silnika liniowego VCM 019-048-02, firmy Moticont.}
		{moticont}
	
Graf \ref{grafiki/vcm_force_current_plot.eps} prezentuje zależność między siłą silnika a położeniem jego cewki. Jak widać największa i stała siła ($ 2 \frac{N}{A} $) występuje dopiero kiedy korpus jest trochę wysunięty i zaczyna lawinowo spadać przy końcowym wysuwie. Jeżeli układ mechaniczny, którego częścią jest silnik, nie zostanie wyposażony w odpowiednią blokadę, korpus może wypaść z cylindra ponieważ silnik nie będzie już w stanie go utrzymać.

	& Producent w karcie katalogowej silnika podaje parametr maksymalnej ciągłej mocy ({\em Max continuous power}) i odpowiada ona zazwyczaj prądowi, który płynie przez cewkę w momencie niepełnego wysunięcia (należy oczywiście wziąć pod uwagę fizyczne obciążenie części ruchomej- masę obudowy lub korpusu, w zależności od tego która z części jest przymocowana, a która jest w ruchu). Przy dużym wysunięciu zależność siły do płynącego przez cewkę prądu drastycznie spada i utrzymując korpus długo w takiej pozycji można w bardzo łatwy sposób spalić uzwojenie silnika. \\
	
\end{easylist}

Przy zastosowaniu odpowiedniej konstrukcji układu napędowego (zachowana równoległość części mechanicznych) i sterowania zapewniającego ochronę przed zbyt dużym prądem płynącym przez cewkę, silnik VCM z czasem nie ulega właściwie żadnej degradacji i długość jego działania jest ograniczona żywotnością części mechanicznych układu.










\subsection{Silnik krokowy}

Silniki krokowe są maszynami elektrycznymi bez komutatorowymi. Zazwyczaj wszystkie uzwojenia silnika są częścią stojana, podczas gdy rotor jest magnesem trwałym lub w przypadku silników VRM ({\em Variable Reluctance Motor}) zębatym blokiem z materiału magnetycznie miękkiego (patrz sekcja \ref{subsec:materialy_fer}). Komutacja (czyli zmiana kierunku/ uzwojenia w którym następuje przepływ prądu elektrycznego) musi być wymuszona poprzez zewnętrzne urządzenie sterujące maszyną. Takie sterowniki są zazwyczaj projektowane w taki sposób, aby zapewnić możliwość utrzymania wirnika w niemal dowolnej pozycji kątowej, oraz obracać nim w obu kierunkach. 

\subsubsection{Rodzaje silników krokowych}

Silniki krokowe ze względu na konstrukcję można podzielić na trzy główne grupy: 
\begin{easylist}
	& silniki o zmiennej reluktancji,
	& z magnesem trwałym,
	& hybrydowe.
\end{easylist}

\insertImg{grafiki/reluctance_moment.png}
		{Powstawanie reluktancyjnego momentu obrotowego.}
		{przepiorkowski}

Reluktancja jest parametrem analogicznym do rezystancji elektrycznej, lecz odniesionym do strumienia magnetycznego. Jest to inaczej ,,rezystancja magnetyczna''.

Obrócenie ruchomego elementu o kąt $ \alpha $ spowoduje że będzie on próbował powrócić do położenia $ \alpha = 0 $, w którym reluktancja obwodu jest najmniejsza. W silnikach VRM przepływ prądu stałego przez uzwojenia powoduje, że zęby wirnika ustawiają się naprzeciw zasilanego uzwojenia.

\insertImg{grafiki/vrm_przekroj.png}
		{Silnik o zmiennej reluktancji (VRM)- Przekrój.}
		{przepiorkowski}

W odróżnieniu od VRM, wirnik silnika krokowego PM ({\em Permanent Magnet}- Magnes trwały) nie posiada zębów. Jest wykonany w postaci walca naprzemiennie namagnesowanego biegunami N i S. Specyficzny rodzaj zębów ma rdzeń stojana. \newline
W zależności od kierunku przepływu prądu w uzwojeniu przyciągane są odpowiednie bieguny wirnika (resunek \ref{grafiki/pm_stepper.png}).

\insertImg{grafiki/pm_stepper.png}
		{Zasada działania silnika PM.}
		{ni}
		
Główną zaletą silnika krokowego PM jest zastosowanie magnesów trwałych w stojanie, dzięki czemu nie ma potrzeby stosowanie szczotek (jak w silnikach DC) oraz jego niska cena. Wadą tego typu maszyny jest relatywnie niski moment obrotowy i brak możliwości rozwijania dużych prędkości obrotowych. \\

Silnik hybrydowy łączy w sobie cechy obu rozwiązań, dzięki czemu zostały poprawione takie parametry jak:
\begin{easylist}
	& moment obrotowy, 
	& maksymalna prędkość obrotowa, 
	& rozdzielczość kroku.
	\\
\end{easylist}

Niestety silniki hybrydowe są około 2-3 razy droższe od silników PM. Wirnik silnika HB ({\em Hybrid Motor}) jest zbudowany z uzębionych nabiegunników i magnesu trwałego, powodującego naprzemienne magnesowanie zębów biegunami N i S. Uzębiony stojan konstrukcją przypomina ten z silnika VRM. 

\insertImgSetSize{grafiki/stepper_budowa.png}
		{70}
		{Konstrukcja dwufazowego silnika krokowego hybrydowego.}
		{enggarage}
		
W silniku HB wirujące pole stojana obraca wirnik z jednego położenia do drugiego na zasadzie jak w silniku VRM. Jest to możliwe dzięki przesunięciu ,,północnej'' i ,,południowej'' części wirnika o pół ząbka. Silniki HB dzięki poprawionym parametrom są obecnie najbardziej popularnymi silnikami krokowymi, pomimo niskiej ceny silników PM.
		
\insertImgSetSize{grafiki/stepper_rotor.png}
		{70}
		{Budowa wirnika silnika krokowego hybrydowego dwufazowego.}
		{edw_sierpien_2002}

\subsubsection{Sterowanie}
\label{sss:sterowanie_krokowy}

W pierwszej kolejności należy rozróżnić sposób zasilania silnika krokowego, który może być unipolarny lub bipolarny (rys. \ref{grafiki/stepper_zasilanie_uni_bi.png}).

\insertImg{grafiki/stepper_zasilanie_uni_bi.png}
		{Sposób sterowania (zasilania) silnika krokowego dwufazowego- z lewej unipolarnie, z prawej bipolarnie.}
		{hackaday}

Przy zasilaniu unipolarnym, dla pojedynczego uzwojenia na zewnątrz silnika wyprowadzone są trzy przewody- Oba końce cewki i odczep w połowie jej długości. Podłączając do odczepu ,,+'' zasilania można sterować połówkami uzwojeń zwierając w odpowiedniej kolejności i na odpowiedni czas końcówki do masy. W zasilaniu bipolarnym, prąd płynie zawsze przez całe uzwojenie (z silnika wyprowadzone są tylko dwa przewody na cewkę). Zaletą wariantu unipolarnego jest możliwość znacznego uproszczenia układu sterowania. Poważną wadą natomiast jest to, że prąd podawany jest zawsze tylko na połowę danego uzwojenia, co ma negatywne przełożenie na moment obrotowy. 

\insertImg{grafiki/stepper_simple.png}
		{Uproszczony model silnika krokowego dwufazowego o kroku 90\degree.}
		{howtomecha}

Oddzielną sprawą jest sposób/ sekwencja podawania impulsów prądowych do uzwojeń silnika. Sterowanie pod tym względem dzieli się na:

\begin{easylist}
	& Falowe, 
	& Pełnokrokowe, 
	& Półkrokowe,
	& Mikrokrokowe.
	\\
\end{easylist}

Do pomocy w wyjaśnieniu poszczególnych typów sterowania posłuży rysunek \ref{grafiki/stepper_simple.png}, prezentujący uproszczony silnik krokowy. Maszyna składa się z wirnika zbudowanego jak na rys. \ref{grafiki/stepper_rotor.png} tyle że z ilością zębów (na wirniku i stojanie) pozwalającą osiągnąć pełen obrót w czterech pełnych komutacjach. Silnik posiada 2 uzwojenia A i B nawinięte w taki sposób, że pary A/A' i B/B' są ze sobą połączone (uzwojenie podzielone na dwie sekcje). Przykładając napięcie do cewki A, w zależności od polaryzacji tego napięcia, prąd będzie wpływał przez uzwojenie A i wypływał z uzwojenia A' lub na odwrót. Tak samo ma się sytuacja w przypadku cewki B.

\tocLessLeftNorm{Sterowanie falowe}

W przypadku braku zasilania, wirnik silnika ustawi się w jednej z czterech spoczynkowych pozycji (0\degree, 90\degree, 180\degree lub 270\degree). Dzieje się tak dlatego, że namagnesowany wirnik usiłuje zająć taką pozycję, aby reluktancja obwodu magnetycznego była jak najmniejsza, tzn. aby namagnesowane zęby wirnika były ustawione w jednej linii z dowolną parą zębów stojana. Znikomy moment obrotowy powodujący takie pozycjonowanie to inaczej moment bezprądowy (spoczynkowy). 

\insertImgSetSize{grafiki/stepper_sterowanie_falowe.eps}
		{100}
		{Sekwencja zasilania uzwojeń przy sterowaniu falowym (obrót wirnika z rys. \ref{grafiki/stepper_simple.png} według wskazówek zegara).}
		{oprWlasne}

Rysunek \ref{grafiki/stepper_sterowanie_falowe.eps} prezentuje sekwencję dzięki której wirnik obróci się o trzy pełne obroty w prawo przy zastosowaniu sterowania falowego. Na osi odciętych znajdują się kolejne chwile w czasie, a na osi rzędnych wartość potencjału przyłożonego do danego zacisku (wartości 1 i 0 prezentujące stan wysoki lub niski potencjału). Aby wirnik obracał się w przeciwnym kierunku, sekwencję należy odtworzyć w odwrotnej kolejności. Jak widać na przebiegach, w danym momencie w czasie prąd znajduje się zawsze tylko w jednym z uzwojeń, co nie umożliwia uzyskania maksymalnego dla danej maszyny momentu obrotowego. Pomimo że sterowanie falowe nie zapewnia uzyskania optymalnego momentu obrotowego, to przy zastosowaniu w niektórych aplikacjach jest wystarczające, a ponadto jest jest najprostsze do zaimplementowania.

\clearpage

\tocLessLeftNorm{Sterowanie pełnokrokowe}

W trybie sterowania pełnokrokowego prąd płynie przez oba uzwojenia jednocześnie. Nie ma to wpływu na prędkość obrotową wirnika, ale ma znaczący wpływ na moment obrotowy, który jest dwa razy większy niż w trybie falowym. Pełny obrót wirnika dalej jest osiągany w czterech "ruchach", ale absolutna pozycja kątowa w każdym z nich jest względem trybu falowego przesunięta o 45\degree.

\insertImgSetSize{grafiki/stepper_sterowanie_krokowe.eps}
		{100}
		{Sekwencja zasilania uzwojeń przy sterowaniu pełnokrokowym (obrót wirnika z rys. \ref{grafiki/stepper_simple.png} według wskazówek zegara).}
		{oprWlasne}
		
Rysunek \ref{grafiki/stepper_sterowanie_krokowe.eps} przedstawia przebiegi dla 3 obrotów, przy sterowaniu pełnokrokowym. Kolejne pozycje wirnika są następujące: 45\degree w pierwszym interwale czasowym, 135\degree w drugim, 225\degree w kolejnym i 315\degree w ostatnim (dla danego obrotu). Potem cykl się powtarza. Tryb pełnokrokowy pozwala osiągnąć optymalny dla danego silnika moment obrotowy, ale niestety rozdzielczość obrotu jest ograniczona do ilości kroków na obrót danego modelu maszyny. Typowe wartości kroków na obrót w silnikach krokowych hybrydowych to 100 (krok o rozdzielczości 3.6\degree), 200 (1.8\degree) i 400 (0.9\degree). Aby uzyskać większą rozdzielczość kroku, należy zastosować sterowanie półkrokowe lub mikrokrokowe.

\clearpage

\tocLessLeftNorm{Sterowanie półkrokowe}

Tryb półkrokowy jest połączeniem sterowania falowego i pełnokrokowego. Rozdzielczość obrotowa względem poprzednich trybów jest dwa razy większa, dlatego że uzwojenia naprzemiennie zasilane są pojedynczo/ oba jednocześnie.

\insertImgSetSize{grafiki/stepper_sterowanie_polkrokowe.eps}
		{100}
		{Sekwencja zasilania uzwojeń przy sterowaniu półkrokowym (obrót wirnika z rys. \ref{grafiki/stepper_simple.png} według wskazówek zegara).}
		{oprWlasne}
		
Wadą takiego rozwiązania jest to że w co drugim kroku moment obrotowy będzie mniejszy. Może to spowodować szarpanie wirnika przy dużych prędkościach obrotowych. Przy niskich prędkościach i obciążeniach natomiast, efekt ten rzadko występuje i można uzyskać polepszoną płynność ruchu, w porównaniu do poprzednich trybów. Aby uniknąć efektu zmiennego momentu, można zastosować sterowanie półkrokowe ze zmiennym prądem (rys. \ref{grafiki/stepper_sterowanie_polkrokowe_zmienne.eps}). Na rysunku widać że w momentach kiedy różnica potencjałów występuje tylko na jednym z uzwojeń, jest ona dwa razy większa niż w przypadku kiedy zasilane są obie cewki. Prąd płynący w uzwojeniach jest proporcjonalny do przyłożonego do nich napięcia, dlatego moment się wyrównuje i silnik nie jest narażony na szarpanie. Rozwiązanie te ma jednak wady- Zakładając że dla momentu w czasie, w którym przez oba uzwojenia płynie prąd, wartość tego prądu jest maksymalna (nominalna) dla danego modelu silnika w każdym z uzwojeń, oznacza to że w chwilach kiedy kiedy napięcie jest podawane na tylko jedno uzwojenie, prąd w nim płynący jest większy od nominalnego (tak jak na rys. \ref{grafiki/stepper_sterowanie_polkrokowe_zmienne.eps}). To naraża silnik na nadmierne grzanie, a nawet na uszkodzenie. Z drugiej strony można założyć że nominalny prąd ma płynąć w momencie zasilania jednego uzwojenia, a mniejszy przy zasilaniu dwóch. Ten tryb jest bezpieczny, ale optymalny moment obrotowy nie jest uzyskiwany.

\insertImgSetSize{grafiki/stepper_sterowanie_polkrokowe_zmienne.eps}
		{100}
		{Sekwencja zasilania uzwojeń przy sterowaniu półkrokowym ze zmiennym prądem (obrót wirnika z rys. \ref{grafiki/stepper_simple.png} według wskazówek zegara).}
		{oprWlasne}
		
\tocLessLeftNorm{Sterowanie mikrokrokowe}

Największą rozdzielczość ruchu (teoretycznie nieograniczoną, a praktycznie ograniczoną możliwościami układu sterującego) można uzyskać przy zastosowaniu trybu mikrokrokowego. Aby obrócić wirnik z rys. \ref{grafiki/stepper_simple.png} w najpłynniejszy możliwy sposób, kształt prądu podawany na uzwojenia powinien być naturalny dla danej konstrukcji silnika, czy też dla ruchu który wykonuje element ruchomy- W tym wypadku jest to ruch obrotowy. 

\begin{equation} \label{eq:stepper2}
	\sin^2\theta + \cos^2\theta = 1
\end{equation}

Równanie \ref{eq:stepper2} prezentuje wzór na jedynkę trygonometryczną. Można go użyć, aby uzyskać dowolną pozycję wirnika pomiędzy kolejnymi naturalnymi ułożeniami wirnika. W tym momencie należy rozróżnić dwa rodzaje pozycji kątowych wirnika silnika krokowego- magnetyczną oraz fizyczną. W przypadku silnika z rys. \ref{grafiki/stepper_simple.png} obie pozycje będą się zawsze pokrywać, ze względu na to że potrzeba jedynie czterech pełnych kroków aby uzyskać pełen obrót- kąt ,,magnetyczny'' 45\degree odpowiada obrotowi wirnika o 45\degree. Dla rzeczywistej sytuacji, w której silnik posiada np. 200 kroków na obrót, jeden stopień magnetyczny odpowiadałby już jedynie $ \frac{1}{50} $ stopnia w obrocie. Innymi słowy potrzeba 50 obrotów magnetycznych na 1 obrót fizyczny wirnika.

\insertImg{grafiki/circle_sine.eps}
		{Ćwiartki okręgu i funkcji $ \sin x $ w kartezjańskim układzie współrzędnych}
		{wikimedia}
		
Tak więc, podając na uzwojenia silnika prąd o kształcie funkcji {\it sinus} (na pierwsze uzwojenie) i {\it cosinus} (na drugie), suma ich kwadratów powinna zawsze wynosić jeden. 

\insertImgSetSize{grafiki/stepper_sin_cos.eps}
		{85}
		{Kształt prądów podawanych na uzwojenia A i B w celu uzyskania optymalnej płynności obrotu w silniku krokowym hybrydowym dwufazowym.}
		{oprWlasne}
		
Prezentowane przebiegi z rys. \ref{grafiki/stepper_sin_cos.eps} obrócą wał silnika o 360\degree magnetycznych. Maksymalna amplituda po znormalizowaniu reprezentuje wartość natężenia nominalnego prądu dla danego modelu silnika. Zazwyczaj silniki krokowe są sterowane z przekształtników, które same są zasilane napięciem stałym. Rozdzielczości uzyskiwane w typowych sterownikach wynoszą od 8 do 32 mikrokroków na pełen krok.  \\

\insertImgSetSize{grafiki/stepper_sin_cos_8.png}
		{70}
		{Kształt prądów podawanych na uzwojenia A i B przy sterowaniu mikrokrokowym z rozdzielczością $ \frac{1}{8} $ kroku.}
		{lab_boni_holub}
		
Tabela \ref{tab:stepper_sterowanie} prezentuje podsumowanie omówionych trybów sterowania silnikami krokowymi dla zasilania bipolarnego i unipolarnego.

\insertTab{|c|c|c|c|c|c|c|c|c|}
{%
\hline
\makecell{Rodzaj\\sterowania} & \multicolumn{2}{|c|}{Falowe} & \multicolumn{2}{|c|}{Pełnokrokowe} & \multicolumn{2}{|c|}{Półkrokowe} & \multicolumn{2}{|c|}{Mikrokrokowe} \\ 
\cline{2-9} & Bi & Uni & Bi & Uni & Bi & Uni & Bi & Uni \\ \hline
\makecell{Moment\\obrotowy} & 50\% & 25\% & 100\% & 50\% & \makecell{50-\\100}\% & \makecell{25-\\50}\% & 50\% & 25\% \\ \hline
\makecell{Komplikacja\\sterownika} & średnia & mała & średnia & mała & duża & duża & duża & duża \\ \hline
\makecell{Płynność\\ruchu} & mała & mała & mała & mała & duża & duża & duża & duża \\ \hline
\makecell{Rozdzielczość\\kroku} & mała & mała & mała & mała & średnia & średnia & duża & duża \\ \hline
}
{Cechy poszczególnych trybów sterowania silnikiem krokowym}
{przepiorkowski}
{tab:stepper_sterowanie}

\subsubsection{Zalety i wady silników krokowych}

Możliwość sterowania silnikiem krokowym za pomocą impulsów jest dużą zaletą- Przy odpowiednio dobranym układzie napędowym, do pracy w ograniczonym obszarze dokładności silnik krokowy nie potrzebuje żadnego sprzężenia zwrotnego (praca w otwartej pętli sprzężenia zwrotnego). Teoretycznie pozycja wirnika jest znana po każdej komutacji wykonanej przez układ sterujący. Niestety, wirnik i jego obciążenie posiadają pewną bezwładność co sprawia że przed osiągnięciem ustalonej pozycji występują drgania. \\
		
W przypadkach kiedy wymagana jest duża precyzja sterowania, tak duże drgania są bardzo niekorzystne. Wybrane sposoby eliminacji tych problemów to:

\begin{easylist}
	& Tłumienie mechaniczne poprzez zwiększenie momentu tarcia.
	& Sprzęgła elastyczne i przekładnie paskowe zastosowane w układzie napędowym, w celu redukcji przenoszonych drgań do obciążenia.
	& Wybór pracy półkrokowej lub mikrokrokowej- Im bardziej przebieg prądu jest zbliżony do sinusoidalnego, tym ciszej i stabilniej pracuje silnik krokowy.
	& Zastosowanie silnika krokowego o mniejszym skoku.
	\\
\end{easylist} 

\insertImgSetSize{grafiki/stepper_drgania.png}
		{60}
		{Wykres czasowy ilusytujący położenie wirnika podczas wykonywania skoku}
		{edw_sierpien_2002}

Należy pamiętać o wielkości i częstotliwości oscylacji decyduje głównie masa wirnika oraz wielkość i rodzaj obciążenia.





\subsection{Silniki prądu stałego (DC)}

\subsubsection{Budowa}

\subsubsection{Zasada działania}

\subsubsection{Sterowanie}

\clearpage


\subsection{Silnik bezszczotkowy z wirującym magnesem (BLDC)}

Konstrukcja silnika BLDC ({\it Brushless DC Motor}) jest odwrotnością silnika komutatorowego z magnesem trwałym: uzwojenia znajdują się w stojanie, a wirnik jest wykonany z odpowiednio ukształtowanego magnesu. Maszyny te są zaliczane do grupy silników prądu stałego, ale w praktyce pradem stałym zasilany jest jedynie falownik sterujący silnikiem. 

\insertImgSetSize{grafiki/bldc_przekroj.jpg}
		{50}
		{Przekrój typowego silnika PM BLDC: a) z wirnikiem wewnętrznym, b) z wirnikiem zewnętrznym, 1- jarzmo stojana (twornika), 2- zęby stojana, 3- żłobek z uzwojeniami, 4- wirnik, 5- magnes trwały}
		{krykowski}
		
Silniki BLDC są konstrukcyjnie bardzo podobne do silników krokowych- można przyjąć że silniki BLDC są pewną podgrupą silników krokowych o kroku 180\degree (silniki 2-fazowe) lub 120\degree (3-fazowe). Kryterium odróżniające silniki krokowe od BLDC może być takie, że w większości aplikacji dla poprawnej silnika BLDC jego sterownik musi dysponować sygnałem zwrotnym o aktualnym położeniu kątowym wirnika (np. z czujników Halla). Silniki krokowe natomiast mogą pracować w otwartej pętli sprzężenia zwrotnego.

\insertImgSetSize{grafiki/bldc_gwiazda.jpg}
		{50}
		{Konfiguracja uzwojeń trójfazowego silnika BLDC w gwiazdę}
		{przepiorkowski}
		
Uzwojenia silnika BLDC dla maszyny 3-fazowej są układane w gwiazdę (rys. \ref{grafiki/bldc_gwiazda.jpg}). Prąd z układu sterującego jest doprowadzany do wyprowadzeń A1, B1 i C1 podczas gdy pozostałe końce uzwojeń (A2, B2, C2) są ze sobą zwarte. Oznacza to że w tym rodzaju maszyny, prąd przepływa zawsze przez co najmniej dwa uzwojenia, w jednym z dwóch możliwych kierunków (np. kiedy wyższy potencjał jest na C1, a niższy na A1 wtedy od C1 do A1). Rysunek \ref{grafiki/bldc_sterowanie.eps} przedstawia tą zależność.

\insertImgSetSize{grafiki/bldc_sterowanie.eps}
		{100}
		{Uproszczone przebiegi prądów w uzwojeniach silnika 3-fazowego sterowanego bipolarnie}
		{oprWlasne}
		
Zmiany momentu obrotowego w funkcji kąta obrotu są w silniku trójfazowym niewielkie i można je jeszcze bardziej zredukować stosując impulsy prądowe o kształcie trapezowym. 
Tak jak zaprezentowano na rys. \ref{grafiki/bldc_przekroj.jpg} silniki bezszczotkowe mogą mieć różną konstrukcję wirnika. W przypadku budowy klasycznej wirnik w kształcie walca jest umieszczony wewnątrz stojana z uzwojeniami. W silniku z wirnikiem zewnętrznym magnetyczny wirnik ma kształt kubka i obraca się wokół nieruchomego stojana. \\

Do głównych zalet silników BLDC należy zaliczyć to że zjawisko rezonansu mechanicznego jest o wiele mniejsze niż w przypadku silników krokowych, oraz mogą one uzyskiwać prędkości obrotowe o wiele większe te pierwsze (powyżej 10000 RPM). Ponadto, przy tej samej wartości momentu obrotowego silnik BLDC może być mniejszy i lżejszy od krokowego. Z racji tego że BLDC nie posiada komutatora, a co za tym idzie szczotek, jego żywotność jest bardzo długa i jest praktycznie ograniczona żywotnością mechaniki układu napędowego. Największą niedogodnością jest natomiast to, że z silnikiem bezszczotkowym z wirującym magnesem nie da się w praktyce pracować w otwartej pętli sprzężenia zwrotnego.

\clearpage



% Opis sposobów sterowania silnikami
\section{Zasilanie silników}

Z praktycznego punktu widzenia najwygodniej było by, aby wyprowadzenia silnika były zasilane z idealnych źródeł prądowych. Tym sposobem, możliwe było by wymuszenie przepływ prądu w uzwojeniu o dowolnym natężeniu bez większego problemu. W praktyce niestety zbudowanie takiego falownika było by bardzo trudne, w szczególności jeśli miałby on pracować z silnikami dużej mocy. Dlatego właśnie do sterowania silnikami stosuje się najczęściej metodę modulacji szerokości impulsów napięciowych ({\it Pulse width modulation}). Dzięki temu, uzwojenia maszyny są zasilane ze źródła napięciowego (które są tanie w budowie, np. wszelkiego rodzaju zasilacze lub baterie), do którego dostęp jest włączany i wyłączany w odpowiednim czasie i na określony czas.

\insertImgSetSize{grafiki/mostek_h_pn.eps}
	{80}
	{Mostek typu ,,H'' zbudowany z czterech tranzystorów polowych z izolowaną bramką ({\it Mosfet})}
	{oprWlasne}

Kwestią pozostaje dobranie odpowiednich ,,przełączników'' które doprowadzały by i odcinały napięcie od uzwojenia. Rys. \ref{grafiki/mostek_h_pn.eps} prezentuje przykładowy mostek typu H, zbudowany z czterech tranzystorów polowych. Tak zbudowany układ może posłużyć do wymuszenia przepływu prądu w uzwojeniu silnika (które symbolizuje gałąź z rezystorem R1 i cewką L1) w dowolnym kierunku, w zależności od tego które z tranzystorów są aktualnie otwarte/ zamknięte. Jedna połowa mostka (półmostek) tworzy układ {\it Push-Pull}, który może w punkcie połączenie dwóch tranzystorów (V1 i V2) zapewnić potenciał o wartości przybliżonej do VCC lub GND. Nie zawsze jako klucze są wykorzystywane tranzystory typu Mosfet, w zależności od aplikacji są to też inne półprzewodniki takie jak tranzystory bipolarne, triaki czy tranzystory IGBT ({\it Insulated-gate bipolar transistor}). W komercyjnych sterownikach małej mocy jednak, najczęściej są to właśnie tranzystory polowe, z tą różnicą że górne tranzystory mostka są także tranzystorami typu N ze względu na ich lepsze względem typu P właściwości (Mobilność elektronów jest większa niż dziur, a to te pierwsze są nośnikami ładunków w tranzystorze typu N \cite{tietze}). W takim wypadku jednak do wysterowania górnych tranzystorów potrzebne są dodatkowe układy elektroniczne, lub dodatkowe źródło zasilania, ze względu na to że napięcie na bramce musi być wyższe od napięcia na źródle tranzystora, w celu jego pełnego wysterowania. Diody połączone równolegle do tranzystorów są niezbędne- w przypadku ich braku prąd uzwojenia usiłowałby zaniknąć w czasie równym 0, co jest niemożliwe, ale napięcie samoindukcji osiągneło by wartość od kilkuset woltów do kliku kV. Wtedy energia zgromadzona na indukcyjności może rozładować się na dwa sposoby: 

\begin{easylist}
	& nastąpi przebicie tranzystora lub izolacji uzwojenia z wyładowaniem iskrowym,
	& pojemność pasożytnicza uzwojenia stworzy obwód rezonansowy LC, energia zostanie wypromieniowana w postaci fali elektromagnetycznej i częściowo zamieniona w ciepło poprzez prądy wirowe zaindukowane w rdzeniu \cite{przepiorkowski}.
	\\
\end{easylist}

Dzięki indukcyjnej naturze uzwojeń silników, kluczując napięcie zasilania (w określonych interwałach czasowych) możliwe jest wymuszenie przepływu prądu w uzwojeniu o wartości bliskiej z góry założonej w sterowniku, bez względu na to że maszyna jest zasilana ze źródła napięciowego.

\insertImgSetSize{grafiki/indukcyjnosc.eps}
	{90}
	{Przebieg prądu po podaniu impulsu napięciowego (w momencie $ t = 0 s $) o wartości $ U = 24 V $ na uzwojenie o rezystancji $ R = 4.1 \Omega $ i indukcyjności $ L = 1.9 mH $}
	{oprWlasne}
	
Czas w którym prąd w uzwojeniu osiągnie swoją maksymalną wartość przedstawia stała czasowa cewki $ \tau $:

\begin{equation} \label{eq:ster1}
	\tau = \frac{L}{R}
\end{equation}

Gdzie $ L $ i $ R $ to indukcyjność i rezystancja uzwojenia. Aby obliczyć prąd uzwojenia w dowolnym momencie $ t $ należy użyć wzoru:

\begin{equation} \label{eq:ster2}
	I_L = \frac{U}{R} (1 - \exp^\frac{-t}{\tau})
\end{equation}

Spadek napięcia na cewce podczas ładowania i rozładowywania cewki wyrażany jest wzorem:

\begin{equation} \label{eq:ster3}
	U_L = - L \frac{dI_L}{dt}
\end{equation}


\subsection{Tryby pracy}

Obserwując kierunek przepływu energii w maszynie elektrycznej, można na tej podstawie wyróżnić jej dwa tryby pracy: silnikową oraz generatorową. W stanach ustalonych układu napędowego zjawiska te można zobrazować za pomocą charakterystyk mechanicznych (rys. \ref{grafiki/ch_statyczne.eps}), pokazujących zależność momentu $ M $ od prędkości $ \Omega $. Kształt charakterystyk zależy przede wszystkim od struktury i parametrów układów regulacyjnych i pomiarowych, a w pewnym zakresie również od charakterystyk przekształtnika oraz od typu i parametrów silnika \cite{zawirski}.

\insertImgSetSize{grafiki/ch_statyczne.eps}
	{100}
	{Charakterystyki statyczne silnika: $ \Omega $ - prędkość obrotowa, $ M $ - moment obrotowy}
	{oprWlasne}
	
Moc na wale maszyny jest wyrażana wzorem:

\begin{equation} \label{eq:chs1}
	P = \Omega M
\end{equation}

Kiedy jest dodatnia (kwadrant I i III) oznacza to że maszyna pobiera ją ze źródła napięcia i pracuje jako silnik napędzając element będący obciążeniem. W przypadku kiedy moc jest ujemna (kwadrant II i IV) silnik hamuje obciążenie pracując jako prądnica (praca generatorowa), która poprzez przekształtnik oddaje energię do źródła \cite{ecn_luty_2012}. \\

W zależności od topologii układu przekształtnikowego i zastosowanych algorytmów sterowania, maszyna elektryczna może mieć możliwość pracy tylko w części kwadrantów, lub we wszystkich . 









\subsection{Sterowanie mostka ,,H''}

Omawiany wcześniej mostek ,,H'' (rys. \ref{grafiki/mostek_h_pn.eps}), który służy do dostarczania prądu do uzwojeń może być sterowany na dwa sposoby-- unipolarny i bipolarny. Oba sposoby mają swoje wady oraz zalety i w związku z tym sterownik, który nie potrafi obsłużyć obu tych trybów może być stosowany tylko w ograniczonych aplikacjach. Ponadto, trybów tych nie należy mylić z trybami zasilania silnika krokowego, które zostały omówione w podrozdziale \ref{sss:sterowanie_krokowy}.

\subsubsection{Modulacja bipolarna}

Sterowaniem bipolarnym nazywa się takie sterowanie, w którym dla kolejnych okresów sygnałów sterujących PWM, wartość napięcia na uzwojeniu zmienia się w zakresie od $ V_d $ (napięcia źródła napięciowego zasilającego silnik) do $ -V_d $. Ten sposób wysterowania mostka przedstawia schemat z rysunku \ref{grafiki/sterowanie_bipolarne.png}. Jako obciążenie został przyjęty silnik DC.

\insertImgSetSize{grafiki/sterowanie_bipolarne.png}
	{80}
	{Silnik DC sterowany bipolarnie}
	{ti_motors}
	
Widoczne na schemacie połączenia i zastosowana bramka NOT jest tylko obrazowym odzwierciedleniem sposobu sterowanie implementowanego programowo w sterowniku. Zakładając sytuację bez obciążenia wirnika, wypełnienie impulsu powyżej 50\% wymusi ruch ,,do przodu'', a poniżej 50\%-- ,,do tyłu''. Taka maszyna jest elektrycznym odpowiednikiem mechanicznej przekładni bezstopniowej CVT ({\it Continuously Variable Transmission}) \cite{ti_motors}. W tego typu maszynie pozycję spoczynkową otrzymuje się poprzez zadawanie sygnału o wartości środkowej. \\

Przy sterowaniu bipolarnym, maszyna pracuje częściowo w czwartym kwadrancie charakterystyki (rys. \ref{grafiki/ch_statyczne.eps}). Tak długo, jak przyłożone do silnika średnie napięcie jest tej samej polaryzacji co SEM generowana w uzwojeniu, i jego amplituda jest wyższa od amplitudy tej SEM, wtedy maszyna pracuje w trybie silnikowym. W przypadku kiedy napięcie zasilania i SEM mają tą samą polaryzację, lecz SEM ma wyższą amplitudę to maszyna pracuje w trybie generatorowym. 

\insertImgSetSize{grafiki/modulacja_bipolarna.png}
	{120}
	{Przebiegi napięć w modulacji bipolarnej}
	{namboodiri_wani}
	
Rys. \ref{grafiki/modulacja_bipolarna.png} prezentuje przebiegu towarzyszące modulacji bipolarnej, zastosowanej w do przykładu mostka z rys. \ref{grafiki/sterowanie_bipolarne.png}. Zadany jest sygnał sinusoidalny $ V_m $, a PWM jest taktowany przebiegiem $ V_{cr} $. $ V_{AG} $ i $ V_{BG} $ to napięcia odpowiednio między potencjałami A i B a masą, natomiast $ V_{AB} $ to spadek na silniku. \\

Zaletą sterowania bipolarnego jest to że z sprzętowego punktu widzenia wymaga bardzo niewiele zasobów. W przypadku najprostszego sterowania wymagany jest tylko jeden sygnał PWM. Oznacza to że nawet najtańszy mikroprocesor może posłużyć do zbudowania sterownika do silnika DC. Ważniejszą zaleta jednak jest to, że używając tylko jednego rezystora pomiarowego ($ R_m $ na rys. \ref{grafiki/sterowanie_bipolarne.png}) istnieje możliwość pomiaru prądu płynącego w uzwojeniu, w dowolnym momencie. Daje to możliwość swobodnego doboru sposobu próbkowania prądu i upraszcza układ sterujący. Najpoważniejszą wadą sterowania bipolarnego jest większa pulsacja prądu w stosunku do metod opisanych w dalszej części pracy.

\subsubsection{Sterowanie unipolarne}
\label{sss:moduni}

Zasadę działania sterowania unipolarnego można zaprezentować bazując na rysunku \ref{grafiki/sterowanie_unipolarne.png}. Przy jego pomocy możliwe jest wysterowanie silnika w taki sposób, aby jego wirnik mógł obracać się w obu kierunkach (kwadranty I i III). Aby pracować w pierwszym kwadrancie, tranzystor Q1 jest włączony, a na bramkę Q4 podawany jest sygnał PWM. Kiedy Q4 jest otwarty, prąd płynie od $ V_d $ przez Q1, dalej przez uzwojenie i przez Q4 do masy poprzez rezystor $ R_m $. Kiedy Q4 się zamyka w uzwojeniu dalej jest prąd ze względu na jego indukcyjność. Zwiększający się potencjał w punkcie połączenia drenu Q4 i źródła Q3 sprawia że dioda tranzystora Q3 zaczyna przewodzić. Z racji tego że Q1 jest cały czas otwarty, obwód prądowy zamknie się właśnie przez ten tranzystor, a nie przez obwód zasilania. Aby obracać wirnikiem w drugą stronę, należy otworzyć Q3 i kluczować Q2. W takim wypadku silnik pracuje w trzecim kwadrancie charakterystyki statycznej (negatywna prędkość i moment. Niezależnie od kierunku przepływu prądu w silniku, napięcie na niego podawane zawsze zmienia się w graniczy od $ V_d $ do GND, a nie jak w schemacie obrazującym sterowanie bipolarne (rys. \ref{grafiki/sterowanie_bipolarne.png}) od $ V_d $ do $ -V_d $. Sprawia to że w takim układzie silnik nie może oddawać energii do zasilania, bo prąd jest ,,zamknięty'' w oczku utworzonym z Silnika, Q1 i Q3. 

\insertImgSetSize{grafiki/sterowanie_unipolarne.png}
	{80}
	{Silnik DC sterowany unipolarnie}
	{ti_motors}


Zaletą prezentowanej topologii jest to że wymagany jest tylko jeden sygnał PWM do sterowania pojedynczym mostkiem (oraz linia wejściowa ustalająca kierunek obrotów). Ponadto, w danym momencie kluczowany jest tylko jeden tranzystor, a więc co za tym idzie straty spowodowane przełączaniem są mniejsze niż w sterowaniu bipolarnym. Modulacja unipolarna wywołuje mniejszą amplitudę tętnień prądu niż modulacja bipolarna przy tej samej częstotliwości modulacji \cite{zawirski}. 

\insertImgSetSize{grafiki/modulacja_unipolarna.png}
	{120}
	{Przebiegi napięć w modulacji unipolarnej}
	{namboodiri_wani}

Bez względu na zastosowaną technikę modulacji, do wszystkich silników zasilanych napięciem stałym można użyć mostka typu ,,H'' w celu realizacji tej modulacji. \\

Konstruowany sterownik korzysta z modulacji unipolarnej, ze względu na to że przynosi o wiele więcej korzyści niż bipolarna. Wadą modulacji unipolarnej jest brak możliwości aktywnego hamowania. Dla urządzenia nie jest to jednak żadnym problemem, ponieważ sposób w jaki uzwojenia są sterowane wyklucza takie sytuacje (więcej informacji na ten temat znajduje się w rozdziale \ref{s:oprogramowanie}.






\clearpage






% Rozwiazania komercyjne
\section{Rozwiązania komercyjne i Open Source}
\label{s:sterowniki}

Na rynku dostępnych jest bardzo wiele rozwiązań pozwalających w szybki sposób zacząć sterować konkretnym rodzajem silnika. Ceny poszczególnych sterowników różnią się od siebie znacząco, w taki sam sposób jak różni się ich funkcjonalność. Są też dostępne projekty bezpłatne ({\it Open Source i Open Hardware}) w rozumieniu konkretnych rozwiązań technologicznych, czy też {\it ,,Know How''}. Innymi słowy schematy elektryczne i mechaniczne, dane {\it CAD'owskie} i oprogramowanie są udostępnione za darmo, a samo urządzenie trzeba fizycznie wykonać samemu. W przypadku tych drugich należy także dokładnie zaznajomić się z licencją na zasadach jakiej są udostępnione, gdyż bardzo często zdarza się że nie można ich stosować bezpłatnie do zastosowań komercyjnych. 

\subsection{Sterowniki jedno-układowe}

Pierwszą grupą sterowników są urządzenia, które dostarczają użytkownikowi jedynie część funkcjonalności związanej ze sterowaniem silnikiem- najczęściej tą która jest najtrudniejsza do zaimplementowania programowo lub sprzętowo na płytce PCB ({\it Printer Circuit Board}) urządzenia. Ich cena w porównaniu do układów zapewniających pełne sterowanie jest wielokrotnie niższa, a ponadto niejednokrotnie są one częścią tych układów. 

\insertImgSetSize{grafiki/smt_packages.png}
	{80}
	{Przykładowe typy obudów w jakich są produkowane układy scalone do sterowania silnikami. Na górze obudowy TSSOP ({\it Thin Shrink Small Outline Package}) i SSOP ({\it Shrink small outline package}) na dole}
	{oprWlasne}

Funkcjonalność danego układu mocno zależy od jego ceny- wraz z jej wzrostem algorytmy zaimplementowane w sterowniku są coraz lepsze, a dodatkowych potrzebnych komponentów elektronicznych do ich poprawnego działania potrzeba coraz mniej (np. końcówki mocy takie jak mostki H są zintegrowane w układzie). Ogólnie rzecz ujmując, urządzenia te stają się coraz bardziej ,,samodzielne''. Ponadto sterowniki do obsługi silników, które wymagają bardziej skomplikowanego sterowania (np. Krokowe lub BLDC) są droższe od tych, które służą do zasilania prostych maszyn (np. silniki komutatorowe DC). 

\subsubsection{AD3950 - Allegro MicroSystems}
\label{ss:allegro}

Układ AD3950 produkowany przez firmę Allegro MicroSystems jest urządzeniem służącym do sterowania silnikami komutatorowymi DC, aczkolwiek nadaje się też do silników VCM. Najwyższe napięcie przy jakim może pracować to 36 V, a chwilowy prąd szczytowy który może dostarczyć do uzwojenia to $ \pm2.8 A $.

\insertImgSetSize{grafiki/A3950_FBD.eps}
	{120}
	{Funkcjonalny schemat blokowy układu A3950}
	{allegro}

Układ posiada w swojej obudowie mostek H oraz ,,logikę'' pozwalającą na sterowanie nim używając pojedynczego sygnału PWM w celu ustalenia prędkości obrotowej/ momentu i kilku lini typu I/O ({\it Input/ Output}), które służą między innymi do zmiany kierunku obrotów (poprzez zmianę kwadrantów pomiędzy pierwszym, a trzecim) i trybu hamowania ({\it Slow Decay} jak w sterowaniu unipolarnym lub {\it Fast Decay} w celu szybkiego hamowania, tak jak przy sterowaniu bipolarnym, tyle że sterownik dodatkowo nie pozwala na wejście silnika w tryb generatorowy zamykając odpowiednie klucze). \\

Urządzenie jak na swoją cenę detaliczną wynoszącą \$3.9 (\cite{digikey}) posiada naprawdę duże możliwości- Zbudowanie w tej cenie podobnego urządzenia z komponentów dyskretnych było by niemożliwe. W prostych aplikacjach w których potrzebne jest sterowanie silnikiem DC niskiej mocy układ sprawdza się znakomicie.

\subsubsection{DRV8825 - Texas Instruments}

Kontroler DRV8825 firmy Texas Instruments jest bardzo popularnym układem zarówno wśród hobbystów jak i profesjonalnych użytkowników. Urządzenie służy do  zasilania pojedynczego silnika krokowego bipolarnego lub dwóch silników DC.

\insertImgSetSize{grafiki/drv8825_FBD.eps}
	{150}
	{Funkcjonalny schemat blokowy układu DRV8825}
	{ti}

Ze względu na dużą ilość wyprowadzeń układ może wydawać się skomplikowany w obsłudze, lecz po wstępnej konfiguracji jego obsługa jest podobna do układu A3950 (sekcja \ref{ss:allegro}). Sterowanie układem polega na podawaniu na wejście STEP impulsów napięciowych, gdzie w zależności od dokonanej za pomocą linii I/O, każdy taki impuls sprawia że silnik robi pełen krok lub mikrokrok (możliwość ustawienia ziarnistości od $ \frac{1}{2} $ do $ \frac{1}{32} $ kroku). Zmieniając częstotliwość podawania impulsów i stany linii MODE0 - MODE2 możliwe jest w miarę płynne sterowanie prędkością obrotową silnika krokowego. W urządzeniach takich jak skanery czy drukarki, w których prędkość obrotowa jest zazwyczaj stała, a wymagana dokładność pozycji wysoka, kontroler DRV8825 sprawdza się bardzo dobrze. Gorzej radzi sobie w aplikacjach w których należy dynamicznie zmieniać tryby pracy silnika (sekcja \ref{sss:sterowanie_krokowy}), a rozdzielczość wymaganego kroku jest bardzo duża, np $ \frac{1}{64} $ kroku. \\

Największa zaleta urządzenia jest taka, że odciąża on sterujący nim układ nadrzędny od rygorystycznego, wykonywanego w czasie rzeczywistym sterowania silnikiem krokowym. Znacząca wada jest natomiast taka, że pomimo że możliwe jest zadanie maksymalnego prądu jakim mają być zasilane uzwojenia, to nie ma możliwości jego bezpośredniego pomiaru przez układ nadrzędny.

\subsubsection{TMCC160 - Trinamic}

Trzecim prezentowanym układem jest sterownik do zasilania silników BLDC i PMSM firmy Trinamic- TMCC160. Poziom integracji tego urządzenia jest o wiele wyższy niż w przypadku dwóch pozostałych chipów (jest tak duży sklasyfikowanie go jako pełnowymiarowego pewnie nie było by błędem). W jego strukturze znajduje się potężny mikrokontroler typu Cortex® M4 z rodziny ARM, dzięki czemu komunikacja z może odbywać się poprzez jeden z dostępnych protokołów szeregowych. Na pokładzie jest także przetwornica DC-DC więc dodatkowe zewnetrzne przetworniki napięcia także nie są wymagane do zasilania logiki układu. 

\insertImgSetSize{grafiki/TMCC160_SBD.eps}
	{80}
	{Ogólny schemat blokowy układu TMCC160}
	{trinamic}

Urządzenie jest jedynie tzw. {\it Gate Driver'em}, co oznacza że do pracy potrzebuje zewnętrznych końcówek mocy (co może być zaletą, gdyż dzięki temu można go dostosować do aplikacji niskiej i dużej mocy). TMCC160 potrafi nawet współpracować z czujnikami Hall'a i enkoderami inkrementalnymi- ma zaimplementowane algorytmy sterowania w zamkniętej pętli sprzężenia zwrotnego. \\

Po dodaniu kilku elementów pasywnych i złącz do płytki PCB z TMCC160, użytkownik otrzymuje funkcjonalny serwo-mechanizm, który może być sterowany zarówno z mikrokontrolera o bardzo niskich zasobach (wysoki poziom integracji) jak i z komputera, dzięki mnogości dostępnych interfejsów szeregowych (np. RS232 lub CAN).

\subsection{Sterowniki pełnowymiarowe}

Druga grupa sterowników to urządzenia samodzielnie, które do pracy potrzebują jedynie podłączenia do źródła zasilania, sygnału sterującego z układu wyższego poziomu (najczęściej interfejs szeregowy) jeśli komunikacja nie jest bezprzewodowa i samego silnika. Jeśli układ pracuje jako serwomechanizm to wymagane jest jeszcze podłączenie urządzenia dostarczającego sygnał zwrotny, np enkodera. Istnieje wiele rodzajów kontrolerów tego typu, od takich do sterowania micro silnikami, po potężne falowniki do silników dużej mocy. 

\subsubsection{Simplexmotion 100A - Simplex Motion AB}
\label{sss:simplex}

Urządzenie Simplexmotion 100A jest sterownikiem o najwyższym poziomie integracji, ze względu na to że nawet sterowany silnik znajduje się w obudowie układu.

\insertImgSetSize{grafiki/SIMPLEXMOTION_100A.jpg}
	{100}
	{Inteligentny, wysoce zintegrowany serwo mechanizm Simplexmotion 100A szwedzkiej firmy Simplex Motion. Wymiary: 72 x 64 x 74 mm, waga: 660 gram. Cena: \$353}
	{simplex}

Maszyna znajdująca się w urządzeniu to bezszczotkowy silnik BLDC prądu stałego o mocy 100 W (400 W w ,,piku'') i maksymalnej rozwijanej prędkości obrotowej 6000 rpm. Zamykając w jednej obudowie silnik i układ sterowania, producent zwolnił użytkownika z dokonywania wielu niezbędnych połączeń potrzebnych do poprawnej pracy samego silnika, np. okablowanie do czujników Hall'a potrzebnych do poprawnej pracy BLDC. Wyprowadzone na zewnętrzne zaciski zostały jedynie:

\begin{easylist}
	& interfejs usb do wygodnego podłączenia układu do komputera w celu tuningu parametrów lub samego sterowania,
	& interfejs RS485 działający w standardzie Modbus,
	& wejście na kwadraturowy enkoder inkrementalny,
	& iyprowadzenia ogólnego przeznaczenia (I/O, PWM, wejścia analogowe),
	& zaciski zasilania (12 - 24 V),
	\\
\end{easylist} 

Ponadto producent zapewnia API (.Net, Linux i LabView) dzięki któremu użytkownicy mogą w prosty sposób zaprojektować własną aplikację sterującą tym serwomechanizmem z komputera. Bez wątpienia wysoki poziom integracji Simplemotion 100A jest dużą zaletą w sytuacji kiedy zawarty w nim silnik odpowiada naszej aplikacji i czas na wdrożenie rozwiązania jest niewielki. Cecha ta jednak może wielokrotnie okazać się czynnikiem wykluczającym w danej aplikacji, ze względu na to że silnik nie odpowiada wymaganiom. Dodatkowo, pomimo że wymiary urządzenia są naprawdę kompaktowe (większość obudowy wypełnia silnik) to nie zmieści się ona wszędzie tam, gdzie zmieściłby się sterownik oddzielony od sterowanej maszyny.

\subsubsection{Mechaduino - Tropical Labs}

Kolejny sterownik służy do sterowania silnikami krokowymi i jest projektem Open Source/ Hardware, udostępnionym na zasadach licencji {\it Creative Commons Attribution Share-Alike 4.0 License}, która mówi o tym że projekt może być używany i modyfikowany nawet w celach komercyjnych, o ile użytkownik w swoim produkcie umieści informacje o tym że bazuje na Mechaduino \cite{crcommons}.  

\insertImgSetSize{grafiki/mechaduino.JPG}
	{80}
	{Sterownik do silników krokowych Mechaduino firmy Tropical Labs, kompatybilny sprzętowo i programowo z platformą Arduino Zero}
	{mechaduino}
	
Projekt bazuje na sprzęcie z platformy Arduino Zero, co oznacza że firmware pracujący na sterującym mikrokontrolerze ATSAMD21G18 (Rdzeń ARM Cortex® M0+) może być napisany w środowisku programistycznym Arduino. Oprogramowanie takie pisze się szybciej niż w C czy assemblerze dzięki wysokopoziomowemu framework'owi. Wydajność takiego kodu jednak pozostawia wiele do życzenia.

\insertImgSetSize{grafiki/arduino_zero.jpg}
	{70}
	{Platforma Arduino Zero firmy Arduino}
	{arduino_zero}
	
Podobnie jak sterownik przedstawiony w poprzedniej sekcji (\ref{sss:simplex}), to urządzenie także posiada płytkę sterownika połączoną z silnikiem (Pasuje do silników w standardzie NEMA17). PCB jest montowane za silnikiem, używając śrub na tulejkach dystansowych. Ze względu na otwartość konstrukcji jednak, nie jest problemem aby płytkę zdjąć i umieścić w innym miejscu, o ile w danej aplikacji sterowanie w otwartej pętli sprzężenia zwrotnego jest wystarczające. 

\insertImgSetSize{grafiki/mechaduino_encoder.eps}
	{70}
	{Płytka Mechaduino z oznaczonym zielonym polem miejscem, w którym znajduje się magnetyczny, absolutny enkoder obrotowy AS5047 firmy AMS}
	{mechaduino}
	
W przypadku kiedy praca w zamkniętej pętli jest niezbędna, płytka musi pozostać na silniku ze względu na to że na jej dolnej warstwie znajduje się enkoder magnetyczny (rys. \ref{grafiki/mechaduino_encoder.eps}), który współpracuje z magnesem przyklejonym na końcu wału silnika. W tym trybie, dokładność pozycjonowania wału silnika wynosi 0.1\degree po dokonaniu procedury kalibracji enkodera. \\

Kompaktowa końcówka mocy znajdująca się na PCB to układ A4954 firmy Allegro Microsystems. Są to zintegrowane w układzie scalonym dwa pełne mostki H, które mogą dostarczyć do silnika prąd o natężeniu szczytowym 2 A. Komunikacja z urządzeniem odbywa się poprzez port szeregowy USB lub $ I^2C $ ({\it Inter-Integrated Circuit}). \\

Niezaprzeczalną zaletą urządzenia jest otwartość kodu, dzięki czemu można dostosować sterownik do potrzeb swoich aplikacji na bardzo niskim poziomie. Mechaduino jest jednak projektem hobbystycznym, którego oprogramowanie i hardware  na dzień dzisiejszy nie są w pełni dopracowane, przez co nie nadaje się jeszcze do profesjonalnego zastosowania.

\subsubsection{Sterowniki firmy Technosoft}

Technosoft jest wiodącą firmą w zakresie opracowywania specjalizowanych rozwiązań i produkcji systemów napędowych. Ilość sterowników dostępnych w ofercie firmy jest bardzo duża, dlatego w tej sekcji zostanie opisana tylko jedna rodzina, która była inspiracją do stworzenia prezentowanego dalej w pracy sterownika do silników.

\insertImgSetSize{grafiki/PIM2403.jpg}
	{80}
	{Przykładowy sterownik z rodziny {\it Plug-in Intelligent Servo Drives}, PIM2403. Firma Technosoft}
	{technosoft}
	
Płytki {\it Plug-in} są urządzeniami o wysokim poziomie integracji, ale inaczej niż poprzednio przedstawione sterowniki nie posiadają zintegrowanego silnika. Na pokładzie znajduje się jednostka centrala- szybki mikrokontroler zmiennoprzecinkowy, który odpawiada za całe sterowanie. Na PCB nie ma żadnych sterowników pośredniczących, wszystkie algorytmy związane ze sterowaniem silnikiem są wykonywane w mikrokontrolerze. Rodzaje silników obsługiwanych przez kontroler to:

\begin{easylist}
	& Krokowy,
	& BLDC,
	& PMSM,
	& DC, 
	& VCM.
	\\
\end{easylist}

Sama płytka sterownika nie wystarczy jednak do rozpoczęcia pracy, bo znajdują się na niej tylko pośrednie peryferia i brakuje odpowiednich wyprowadzeń. 

\insertImgSetSize{grafiki/PIM2403_wboard.jpg}
	{80}
	{Kontroler PIM2403 zamontowany na płytce bazowej MS02-06-31 firmy Mechatronic Systems}
	{oprWlasne}

Rysunek \ref{grafiki/PIM2403_wboard.jpg} prezentuje sterownik {\it Plug-in} zamontowany na przykładowej płytce bazowej, która umożliwia komunikację oraz podłączenie do sterownika jako całości wielu końcowych peryferiów potrzebnych do działania. Poprzez złącza można podłączyć:

\begin{easylist}
	& zasilanie, 
	& silnik,
	& enkoder inkrementalny,
	& czujnik Hall'a działający jako krańcówka,
	& interfejsy CAN i RS232,
	\\
\end{easylist}

Jest to przykładowa konfiguracja, która została przystosowana do konkretnej aplikacji- Firma Technosoft udostępnia aplikację PC-ową {\it Easy Motion Studio}, dzięki której możliwe jest nie tylko ustawianie parametrów regulatorów itd., ale nawet zaprogramowanie sterownika w niemal dowolny sposób, używając opracowanego przez firmę języka skryptowego {\it Motion Programming Language}. \\

Rozwiązania oferowane przez tego producenta cechuje możliwość naprawdę szybkiego wdrożenia do projektu. Za tą zaletę trzeba dużo zapłacić, a cena rośnie wraz ze wzrostem możliwości ,,mocowych'' sterowników pod względem zasilania silników. Skala wzrostu jest niestety nieuwarunkowana, bo jedyne co zmienia się na sterownikach to końcówki mocy i czasami rozmiar PCB, a mimo wszystko cena potrafi wzrosnąć nawet o 60 \euro{} (różnica w cenie między sterownikiem 3 i 5-cio amperowym). 

\tocLessLeftNorm{Podsumowanie}

Na rynku jest dostępnych bardzo wiele rozwiązań pozwalających na sterowanie wszelkiego rodzaju silnikami elektrycznymi. Przedstawione przykładowe kontrolery stanowią tylko małą część dostępnych na rynku produktów, ale ukazują formy w jakich buduje się sterowniki. Najtańszymi układami są jedno-chipowe urządzenia, które do poprawnej pracy potrzebują bardzo wielu elementów dodatkowych. Wraz z poziomem integracji cena rośnie, a dodatkowo trzeba dopłacić za możliwość szybkiego wdrożenia produktu do swojej aplikacji i uniwersalność (pod względem wykonywanych programów i obsługi silników). \\

Pomimo wszystkich zalet jakie posiadają zaprezentowane w rozdziale \ref{s:sterowniki} urządzenia, żadne z nich nie jest idealne- najmniejsze układy mają za małe możliwości, a wysoko zintegrowane sterowniki są za drogie (i jak na swoją cenę, także mają za małe możliwości). Z tych właśnie powodów powstaje uniwersalny sterownik do silników, którego dokładny opis znajduje się w dalszych rozdziałach pracy.







\clearpage




% Koncepcja sterownika
\section{Koncepcja projektowanego sterownika}

Istnieje wiele powodów które mogły by warunkować zaprojektowanie kompletnego sterownika do silników, zamiast skorzystania z gotowych rozwiązań (np. prezentowanych w rozdziale \ref{s:sterowniki}).

\insertImgSetSize{grafiki/p10.png}
	{150}
	{Przykładowy produkt firmy \firma{}: Automat P10- maszyna typu {\it Pick and Place} (układająca komponenty elektroniczne na płytkach PCB) z wbudowanym dyspenserem pasty lutowniczej/ kleju.}
	{mechsys}
	
Firma \firma{} zajmująca się dostarczaniem produktów i rozwiązań z dziedziny technologii SMT ({\it Surface Mount Technology}), dla której konstruowany jest sterownik stanowiący temat niniejszej pracy używa w swoich konstrukcjach sterowników Plug-in firmy Technosoft, prezentowanych w sekcji \ref{technosoft}. Postanowiono zastąpić je własnym rozwiązaniem, dlatego że:

\begin{easylist}
	& po dokonaniu rachunku kosztów okazało się, że długoterminowo opracowanie własnego kontrolera będzie bardziej opłacalne niż regularny zakup gotowych sterowników;
	& technologia dostarczana przez Technosoft przestała spełniać wymagania dynamicznie rozwijających się aplikacji w firmie \firma{};
	& tworząc własny sterownik uzyskuje się bardzo duże doświadczenie z dziedziny napędów i sterowania;
	\\
\end{easylist}

Zostaną dalej przedstawione najważniejsze wymagania i warunki brzegowe, które musi spełnić konstruowane urządzenie w celu zastąpienia nim sterownika Technosoftu.

\subsection{Uniwersalność}

\firma{} dostarcza kilka rodzajów produktów, a każdy z nich pracuje z innymi silnikami (Tabela \ref{tab:maszyny_silniki}).

\insertTab{|c|c|c|}
{%
\hline Maszyna & Stosowane silniki & Stosowane sterowniki \\
\hline \makecell{Automaty pick and place \\ i dyspensery}  & krokowe, VCM, DC & Technosoft \\
\hline Drukarki szablonowe & BLDC & \makecell{inne pełnowymiarowe \\ dedykowane} \\
\hline Podajniki automatyczne & DC, krokowe & własne dedykowane \\
\hline Konwejery & DC, BLCD & własne dedykowane \\
\hline
}
{Typ maszyny produkowanej przez \firma{} i używane przez nią silniki wraz ze sterownikami}
{oprWlasne}
{tab:maszyny_silniki}

Na obecnym etapie w części urządzeń stosowane są uniwersalne sterowniki Technosoft'u, które są odpowiednio zaprogramowane do pracy z określonym rodzajem silnika lub własne proste sterowniki, dedykowane pod konkretną maszynę elektryczną (np. silnik komutatorowy DC).
	
Rysunek \ref{grafiki/maszyna_osie.png} prezentuje automat P30 z oznaczonymi osiami. Poprzez pojedynczą oś należy rozumieć jeden układ napędowy w którym zastosowany jest jeden silnik i jeden kontroler sterujący. W tego rodzaju maszyny jest aż sześć układów napędowych, a co za tym idzie sześć zastosowanych sterowników typu Plug-in Technosoft'u.

\insertTab{|c|c|}
{%
\hline Oś & Funkcja \\
\hline Y & Ruch ramienia automatu ,,Prawo- Lewo'' \\
\hline X & Ruch ramienia automatu ,,w przód- w tył'' \\
\hline Z1 & Ruch głowicy układającej ,,Góra- dół'' \\
\hline R1 & \makecell{Obrót głowicy układającej CW i CCW \\ ({\it Clockwise/ Counter clockwise})} \\
\hline Z2 & Ruch głowicy dyspensującej ,,Góra- dół'' \\
\hline R2 & Dyspensja medium z głowicy dyspensującej  \\
\hline
}
{Opis osi w automacie P30}
{oprWlasne}
{tab:maszyny_osie}

Fakt zastosowania tak dużej ilości sterowników znacząco podwyższa koszt jednostkowy wyprodukowania maszyny. Chcąc użyć w napędach poszczególnych osi własnego kontrolera, musi mieć on możliwość sterowania ostatecznie wszystkimi tymi silnikami elektrycznymi, co obecnie sterowniki Technosoft'u. 

\insertImgSetSize{grafiki/maszyna_osie.png}
	{110}
	{Model 3D maszyny P30 typu Pick and Place z wbudowanym dyspenserem z oznaczonymi osiami, a) widok lufy i głowic od przodu maszyny, b) Cała maszyna w rzucie ortogonalnym}
	{mechsys3d}

Dodatkowo, sterownikowi zostało postawione jeszcze jedno wymaganie, którego nie spełnia żaden z prezentowanych kontrolerów dostępnych na rynku: praca z dwoma dowolnymi silnikami równolegle (patrz sekcja \ref{ss:axis_virtualisation}). Ta cecha sprawia, że gotowe rozwiązania pozostają daleko w tyle pod względem funkcjonalnym. Koszt jednostkowy maszyny obniży się nie tylko przez to, że każdy ze sterowników zostanie zastąpiony własnym, ale dodatkowo dlatego, że będzie trzeba zastosować dwa razy mniej kontrolerów. \\

Należy zaznaczyć że krótkoterminowo i w przypadku niewielu wyprodukowanych maszyn, projektowanie własnego kontrolera nie jest opłacalne, gdyż koszt projektu będzie przekraczać koszt zastosowania gotowego rozwiązania. \linebreak Z perspektywy czasu jednak sytuacja się odwraca. Cena za komponenty elektroniczne i wyprodukowanie PCB w zewnętrznej firmie jest niewielka przy zamówieniach w dużych ilościach. Głównym kosztem jest praca wykonana przez projektanta PCB  i programistę w celu opracowania rozwiązania. Zalet tego typu podejścia jest jeszcze więcej i zostaną one w większych szczegółach przytoczone w dalszych rozdziałach i sekcjach.

\subsection{Zgodność protokołu komunikacyjnego}

Budowany sterownik musi zastąpić kontroler Technosoft' zachowując przy tym możliwie wysoki poziom zgodności. W przypadku nie spełnienia tego wymagania, pozostałe moduły (urządzenia) będące częścią całego automatu obarczone byłyby potrzebą aktualizacji ich oprogramowania lub nawet sprzętu, w celu możliwości komunikacji z nowymi sterownikami. 

\insertImgSetSize{grafiki/struktura_sterownikow.eps}
	{150}
	{Struktura zastosowanych sterowników w maszynie P30 obecnie i w przyszłości po pełnej implementacji opracowywanego rozwiązania}
	{oprWlasne}

To wprowadziłoby do przedsięwzięcia bardzo wysokie koszty związane z dodatkowym czasem pracy inżynierów i potrzebą zamówienia nowych przerobionych modułów. Dlatego też warstwa odpowiedzialna za komunikację musi odpowiadać tej zastosowanej przez firmę Technosoft. Specyfikacja do tego protokołu ({\it TML CAN- Technosoft Motion Language}) jest udostępniana przez producenta aby umożliwić użytkownikom ich produktów komunikację z nimi. W tym wypadku zostanie ona zastosowana w celu  ograniczenia czasu potrzebnego na wdrożenie nowego, własnego rozwiązania.

\subsection{Wirtualizacja osi}
\label{ss:axis_virtualisation}

Chcąc zastosować kontroler, który ma możliwość sterowania dwoma silnikami \linebreak w miejsce takiego, który mógł pracować tylko z jednym, trzeba pójść na pewne ustępstwa. Mianowicie, pozostałe urządzenia muszą dalej postrzegać sterownik jako dwa osobne urządzenia. W tym celu kontroler musi przedstawiać się na magistrali komunikacyjnej jako dwa osobne urządzenia, pomimo tego, że posiada tylko jedno oprogramowanie i są do niego podłączone dwa silniki. Taki koncept wydaje się być rozwiązaniem typu {\it work-around} (na około), ale ma także swoje plusy z punktu widzenia architektury oprogramowania kontrolera (patrz rozdział \ref{s:oprogramowanie}). \\

Na chwilę obecną, nowy sterownik jest ciągle w fazie projektowania. W jego pierwotnej formie jest wdrażany do zoptymalizowanej mechanicznie głowicy układającej w automatach Pick and Place. W dalszych rozdziałach i sekcjach, kontroler będzie opisywany w odniesieniu właśnie do tej implementacji. Pomimo tego, że układ może przyjmować nieco zmienioną sprzętową formę w zależności zastosowania \linebreak w różnych aplikacjach, to jego oprogramowanie i on sam dalej pozostaje uniwersalny. W celu optymalizacji zajmowanego przez układ miejsca, hardware może być przystosowany np. tylko do silników krokowych lub tylko DC. Nie oznacza to jednak, że potrzebne są jakiekolwiek zmiany w oprogramowaniu, lub duży wysiłek w przeprojektowaniu płytki PCB. \\ 

Dalsze rozdziały przedstawiają rozwiązania techniczne sprzętowe i programowe w formie szczegółowej.

\clearpage





% Opis hardwaru
\section{Hardware}

Tak jak zostało wspomniane w poprzednim rozdziale, pierwszym zastosowaniem projektowanego sterownika jest nowa, zoptymalizowana pod względem mechanicznym głowica układająca do automatów typu Pick and Place w firmie \firma{}. Na najbliższych stronach znajduje się opis poszczególnych modułów tej głowicy pozwalający zrozumieć jej zasadę działania, ze szczególnym uwzględnieniem części elektrycznej, gdyż to ona stanowi wkład autora w projekt tej głowicy.

\insertImgSetSize{grafiki/glowica.png}
	{140}
	{Prototyp głowicy układającej z zastosowaniem projektowanego sterownika: a) Widok modelu 3D głowicy od boku z wsuniętą pinolą, b) ten sam model z pinolą wysuniętą, c) rzeczywiste zdjęcie głowicy. Wymiary przy wsuniętej pinoli 34 x 38 x 184 mm}
	{mechsys3d}

Rysunek \ref{grafiki/glowica.png} prezentuje sterownik zamontowany w miejscu pracy. Poszczególne elementy urzadzenia to:

\begin{enumerate}
	\item płytka sterownika główna,
	\item płytka sterownika pośrednia zasilająca,
	\item płytka sterownika z układem enkodera dla silnika w osi Z,
	\item płytka sterownika z układem enkodera dla silnika w osi R,
	\item silnik liniowy VCM odpowiedzialny za ruch posuwisto-zwrotny w osi Z,
	\item silnik hybrydowy krokowy odpowiedzialny za obrót w osi R,
	\item pinola głowicy służąca do podnoszenia i odkładania komponentów SMD.
\end{enumerate}

Ze względu na konstrukcję mechaniczną głowicy, w celu wykorzystania jak największej przestrzeni na PCB, kontroler jest ,,rozproszony'' na cztery osobne płytki PCB. Jak widać na rysunku \ref{grafiki/glowica.png} b), elementy 1, 3, 4, 6 i 7 znajdują się w części ruchomej głowicy, a 2 i 5 w części nieruchomej (korpus silnika VCM jest ruchomy, a jego cylinder statyczny- dokładna budowa tego silnika przedstawiona jest na rysunku \ref{grafiki/vcm_budowa.png}). Część ruchoma porusza się po prowadnicy, która zapewnia równoległy względem cylindra ruch korpusu. Zasilanie i sygnały magistrali komunikacyjnej do części ruchomej doprowadzane są elastyczną tasiemką typu {\it Flex} widoczną na zdjęciu głowicy bez obudowy (\ref{grafiki/glowica.png} c)). Silnik krokowy służy do obracania przyssanych do pinoli (podnoszenie przez podciśnienie) komponentów. \\

Kolejne sekcje przedstawiają opis poszczególnych modułów znajdujących się na każdej z płytek PCB kontrolera. Wszystkie schematy ideowe (także ich części) i płytki PCB prezentowane dalej zostały w domyśle wykonane przez autora pracy w programie {\it Altium Designer} dla firmy \firma{}, chyba że zaznaczono inaczej. Niektóre części schematów zostały przeorganizowane w taki sposób, aby zajmowały jak najmniej miejsca w celu czytelnego wyświetlenia w tekście. Kompletny oryginalny schemat ideowy całego urządzenia znajduje się w załącznikach pracy.

\subsection{Płytka zasilająca}
\label{ss:power_board}

Pierwsza z opisywanych płytek na rysunku \ref{grafiki/glowica.png} jest oznaczona numerem 2. Znajduje się na niej zaledwie kilka komponentów- pełni ona funkcję pośredniczącą między płytką główną sterownika a przewodami doprowadzającymi sygnały sterujące i zasilanie.

\insertImgSetSize{grafiki/schemat_plytka_zasilajaca.eps}
	{63}
	{Schemat ideowy płytki zasilającej w głowicy układającej}
	{mechsyspcb}
	
Poprzez przewód do złącza {\it CON1} do głowicy doprowadzane są:

\begin{easylist}
	& zasilanie ze źródła napięcia o amplitudzie 24 V prądu stałego,
	& para różnicowa magistrali CAN,
	& sygnał odbiorczy i nadawczy magistrali RS232.
	\\
\end{easylist} 

Wszystkie powyższe sygnały są następnie przekazywane na złącze {\it CON2}, które poprzez tasiemkę typu {\it Premo-Flex} o pitch'u 0.5 mm (firma Molex) przewodzi je dalej do płytki głównej. Wytrzymałość prądowa pojedynczej żyły takiego złącza to tylko 0.5 A, dlatego sygnały zasilające używają ich aż 8 (po 4 dla zasilania i dla masy).

\insertImgSetSize{grafiki/pcb3d_plytka_zasilajaca.png}
	{70}
	{Widok ortogonalny modelu 3D płytki zasilającej, a) z widokiem warstwy dolnej, b) z widokiem warstwy górnej}
	{mechsyspcb}

Kolejnymi elementami PCB są dwa kondensatory elektrolityczne o największej pojemności jaką tylko udało się dostać w obudowach, które mieszczą się na płytce. Służą po to, aby zwiększyć pojemność widzianą z zacisków zasilania układu, która z kolei jest potrzebną w celu utrzymania stabilnego napięcia w fazach wzmożonego poboru prądu przez urządzenie. Ostatni komponent to zworka zwierająca rezystor terminujący na magistrali CAN. W zależności od topologii magistrali można ją załączyć lub zostawić rozwartą. Terminator został umieszczony dosyć daleko od układu sterownika CAN, ze względu na to że inne miejsce było by zbyt trudno dostępna. Odległość ta jest jednak wystarczająco niska do poprawnego działania układu.

\subsection{Płytka enkodera magnetycznego}
\label{ss:encmag}

Płytka z enkoderem dla osi R jest zamontowana za silnikiem krokowym (numer 4 na rysunku \ref{grafiki/glowica.png}). Jako enkoder magnetyczny został zastosowany układ scalony AS5048A firmy AMS. Rozwiązanie to przypomina te wykorzystane przez Tropical Labs w swoim sterowniku Mechaduino (\ref{ss:mechaduino}). Należy jednak wspomnieć że koncept został zastosowany przed opublikowaniem wpisu o ich sterowniku na stronie \url{hackaday.io} i nie jest na nim wzorowany ani nim inspirowany. AS5048A jest enkoderem absolutnym, który dostarcza informację o aktualnej pozycji kątowej między innymi poprzez interfejs SPI {\it Serial Peripheral Interface}). Używając go można uzyskać pozycję z dokładnością do 0.05\degree{} po zastosowaniu algorytmów uśredniających (\cite{AS5048}). Więcej informacji o zastosowanych w sterowniku algorytmach uśredniających dla tego enkodera jest zawarte w sekcji <dodac referencje>.

\insertImgSetSize{grafiki/AS5048A_diagram.eps}
	{75}
	{Schemat blokowy układu scalonego AS5048A}
	{AS5048}
	
Obracając magnesem w sąsiedztwie układu scalonego, zmienne pole magnetyczne wpływa na matrycę czujników Hall'a która się w nim znajduje. Przetwornik analogowo cyfrowy zawarty w urządzeniu przetwarza analogową informację w formie zmiennego natężenia pola magnetycznego na informację cyfrową o aktualnej pozycji kątowej.

\insertImgSetSize{grafiki/schemat_enkoder_magnetyczny.eps}
	{75}
	{Schemat ideowy płytki z enkoderem magnetycznym dla osi R w głowicy układającej}
	{mechsyspcb}
	
Układ na schemacie ideowym \ref{grafiki/schemat_enkoder_magnetyczny.eps} oznaczony jest symbolem U1. Sygnały magistrali SPI są wyprowadzone na pionowe złącze CON1 o rastrze 1.27 mm, przez które dalej idą na płytkę główną kontrolera. Elementy oznaczone jako FID1 i FID2 to punkty referencyjne dla górnej kamery automatu Pick and Place (tzw. {\it Fidutial'e}). Są one widoczne na płytce PCB z rys. \ref{grafiki/pcb3d_plytka_encmag.png} a), w lewym górnym i prawym dolnym rogu. Ten sam rysunek (obiekt b)) przedstawia płytkę enkodera zamocowaną na elemencie aluminiowym frezowanym, który z kolei jest przykręcony do silnika krokowego. Duży czarny element to mikro zawór doprowadzający podciśnienie, sterowany z głównej płytki kontrolera (opis w sekcji \ref{ss:mainboard}). \\

Oś poprzeczna układu scalonego musi się idealnie pokrywać z osią obracającego się magnesu w celu jak najlepszego odczytu natężenia pola przez enkoder.

\insertImgSetSize{grafiki/pcb3d_plytka_encmag.png}
	{140}
	{Widok płytki enkodera magnetycznego- a) dolna warstwa PCB, na której znajduje się AS5048A wraz z kondensatorami blokującymi, b) rzut ortogonalny płytki zamocowanej na elemencie wspólnym dla niej i silnika}
	{mechsyspcb}
	
Na rysunku \ref{grafiki/przekroj_silnik.png} widać dokładnie co znajduje się w aluminiowej kostce do której zamocowana jest płytka:

\begin{enumerate}
	\item płytka z enkoderem magnetycznym AS5048A,
	\item płytka sterownika główna, połączona przez CON1 (\ref{grafiki/schemat_enkoder_magnetyczny.eps}),
	\item magnes trwały namagnesowany poprzecznie z otworem w środku, zamocowany na końcu wału silnika,
	\item przelotowy wał silnika krokowego,
	\item silnik krokowy hybrydowy,
	\item element aluminiowy (na modelu przezroczysty) łączący wymienione komponenty.
\end{enumerate}
	
\insertImgSetSize{grafiki/przekroj_silnik.png}
	{100}
	{Płytka z enkoderem magnetycznym zamocowana nad obracającym się, przelotowym magnesem trwałym, który z kolei jest zamocowany na przelotowym wale silnika krokowego}
	{mechsys3d}
	
Do elementu łączącego jest także doprowadzone podciśnienie poprzez giętki przewód pneumatyczny- w ten sposób jest możliwe podnoszenie komponentów elektronicznych przez głowicę (przysysając je).
	
\subsection{Płytka enkodera optycznego}
\label{ss:encopt}

Podobnie jak silnik krokowy współpracuje z enkoderem magnetycznym w osi R, tak samo silnik VCM działa z enkoderem optycznym w osi Z i także posiada osobną płytkę PCB dla obwodu tego enkodera. W odróżnieniu do poprzedniej PCB jednak, ta wymaga o wiele więcej komponentów do poprawnej pracy. Pozycja nie jest już podawana absolutnie poprzez magistralę szeregową tylko musi być odczytana jako zmieniający się sygnał analogowy, który następnie poddawany jest interpolacji. Potem trafia do głównej płytki kontrolera w postaci 2 zmieniających się sygnałów zero-jedynkowych (sygnał kwadraturowy). W zależności od tego które zbocze jest wiodące, procesor może określić kierunek przesuwania się enkodera względem liniału. \\

Rysunek \ref{grafiki/enkoder_liniowy_schemat1.eps} prezentuje większość schematu ideowego płytki enkodera. Elementem odczytującym pozycję jest AEDR-8320 (obiekt U2 na schemacie). Jest to enkoder optyczny w małej obudowie SMD, który generuje impulsy świetlne, które odbijając się od przeszkód są przez niego ,,łapane'' przy użyciu matrycy fotodiod.

\insertImgSetSize{grafiki/enkoder_liniowy_schemat1.eps}
	{72}
	{Część schematu ideowego płytki enkodera liniowego osi Z}
	{mechsyspcb}
	
W celu uzyskania najlepszego efektu i rozdzielczości, należy użyć specjalnego paska z mikro wpustkami jako powierzchni odbijającej światło (liniału). Zasada działania jest zobrazowana na rysunku \ref{grafiki/AEDR-8320.pdf}.

\insertImgSetSize{grafiki/AEDR-8320.pdf}
	{60}
	{Schemat blokowy i zasada działania układu scalonego AEDR-8320 firmy Avago Technologies}
	{aedr8320}
	
Sygnały analogowe o kształcie sinusoidy (sygnał wyjściowy A) i cosinusoidy (sygnał wyjściowy B) są podłączone do wejść interpolatora IC-TW (rys \ref{grafiki/ictw2_block_diagram.eps}). Jego zadaniem jest zamiana tych sygnałów na postać cyfrową, która może być odczytana przez mikrokontroler sterujący na głównej płytce urządzenia (połączenie przez złącze CON1). Poziom interpolacji i wzmocnienie ustawia się zwierając odpowiednia wyprowadzenia chip'a do masy lub potencjału zasilającego. Z racji tego że urządzenie jest dopiero prototypem, na PCB enkodera komponenty zostały umieszczony w taki sposób, aby konfigurację można było zmodyfikować przy niewielkim nakładzie pracy- za pomocą zworek, które widoczne są na drugiej części schematu ideowego \ref{grafiki/enkoder_liniowy_schemat1.eps} (rezystorów 0 $ \Omega $).

\insertImgSetSize{grafiki/ictw2_block_diagram.eps}
	{100}
	{Schemat blokowy 8-bitowego interpolatora IC-TW2 firmy IC-HAUS}
	{ictw2}
	
Należy pamiętać o tym aby w danym momencie przylutowany był tylko jeden z rezystorów dla danej lini konfiguracyjnej (zwierający do GND lub 5V). W przeciwnym wypadku połączenie spowoduje zwarcie w układzie.

\insertImgSetSize{grafiki/enkoder_liniowy_schemat2.pdf}
	{40}
	{Część schematu ideowego płytki enkodera liniowego osi Z, prezentująca zworki konfiguracyjne}
	{mechsyspcb}
	
Ostatnim elementem pozostałym do omówienia na schemacie jest U1. SFH9245 to układ elektroniczny działający w podobny sposób do AEDR-8313 (\ref{grafiki/AEDR-8320.pdf}), który jest jednak o wiele prostszy i tańszy. Składa się on z diody emitującej światło podczerwone i z fototranzystora z przerzutnikiem Shmitt'a (\ref{grafiki/sfh9245_schemat_blokowy.pdf}). Światło zamiast odbijać się od liniału, odbija się od wydrukowanego na drukarce laserowej paska, przyklejonego koło liniału. Pasek składa się z dwóch ,,stref''- czarnej i białej.

\insertImgSetSize{grafiki/sfh9245_schemat_blokowy.pdf}
	{50}
	{Schemat blokowy układu scalonego SFH9245 firmy Osram}
	{sfh9245}

Kiedy układ znajduje się nad czarną strefą, światło nie odbija się dobrze i na wyjściu OUT układu jest stan wysoki (wyjście typu {\it Open-Collector}). W strefie białej światło odbija się bardzo dobrze i stan OUT zmienia się na niski. Mechanizm ten jest używany w głowicy do odnajdowania przez silnik liniowy VCM pozycji odniesienia ({\it Homing position}). Jest to niezbędne, dlatego że skonstruowany enkoder jest liniowy, a nie absolutny czyli potrafi dostarczać jedynie informację o zmianie pozycji, a nie ją samą jak poprzedni enkoder (sekcja \ref{ss:encmag}).

\insertImgSetSize{grafiki/enkoder_optyczny_pcb.png}
	{60}
	{Model 3D płytki drukowanej enkodera liniowego silnika VCM. Po lewej warstwa górna, po prawej dolna}
	{mechsyspcb}

Rozdzielczość możliwa do uzyskania w prezentowanym układzie zależy głównie od dwóch czynników: gęstości umiejscowienia rowków na liniale i możliwości interpolatora. W tym przypadku rozdzielczość ogranicza układ IC-TW2, którego maksymalny poziom interpolacji to x64. Przy zastosowaniu liniału 180 LPI ({\it Lines per inch}) daje to rozdzielczość dokładnie 2.20486 mikronów.

\subsection{Płytka główna sterownika}
\label{ss:mainboard}

Czwarta opisywana płytka drukowana była o wiele większym wyzwaniem projektowym, ze względu na dostępne wymiary i dużą ilośc pożądanych elementów. Najważniejsze dane liczbowe płytki:

\begin{easylist}
	& wymiary: 109.1 x 25.6 mm,
	& końcowa grubość laminatu: 1.55 mm,
	& ilość warstw: 4,
	& ilość przelotek: 411,
	& liczba komponentów na warstwie górnej: 145,
	& liczba komponentów na warstwie dolnej: 34.
	\\
\end{easylist} 

Kształt i rozmiar PCB został wybrany tak, aby wykorzystać maksymalnie całą dostępną przestrzeń w obudowie głowicy. Rysunek \ref{grafiki/plytka_glowna_foto.jpg} prezentuję poglądowe zdjęcie płytki niezamontowanej w urządzeniu. 

\insertImgSetSize{grafiki/plytka_glowna_foto.jpg}
	{45}
	{Zdjęcie głównej płytki sterującej głowicy układającej (warstwa górna)}
	{oprWlasne}

Dwie wewnętrzne warstwy płytki to tzw. {\it Plane Layers} i służą głównie do doprowadzania zasilania do układów znajdujących się na płytce. Na warstwie górnej ({\it Component layer}) jest ułożonych ok. 80\% komponentów. Pozostałe znajdują się na warstwie dolnej ({\it Solder side}) tylko dlatego że nie zmieściły się na górze. W momencie kiedy wszystkie komponenty elektroniczne znajdują się tylko po jednej stronie płytki, znacząco upraszcza się jej technologia produkcji. W tym wypadku, aby ,,ułożyć płytkę'' przy pomocy automatu typu Pick and Place, należy po za pastą lutowniczą użyć też kleju (dozowanego np. z dyspensera) aby przykleić komponenty po jednej stronie płytki. Dzięki temu kiedy ta trafi do pieca rozpływowego w celu zlutowania, elementy po dolnej stronie nie pospadają. Taka forma projektowania ma jeszcze jedną zaletę- można użyć jednej z warstw tylko do prowadzenia ścieżek, które inaczej było by trudno doprowadzić w niektóre miejsca. W tym wypadku jednak trzeba było pójść na kompromis; ścieżki (prowadzone lokalnie ,,na krzyż'') i komponenty musiały się znaleźć po obu stronach. Ścieżki sygnałowe nie są prowadzone na wewnętrznych warstwach, aby uniknąć pojemności które mogły by się pojawiać pomiędzy nimi. Warstwy zasilania je od siebie pod względem pojemnościowym izolują.  \\

Elementy THD (przewlekane- {\it Through Hole Device}) zostały ograniczone do minimum w celu optymalizacji czasu montażu. Są to:

\begin{easylist}
	& złącze łączące płytkę główną z płytką enkodera magnetycznego (\ref{ss:encmag}),
	& złącze programatora (\ref{mcusekcja_dodac}),
	& czujnik podciśnienia (\ref{czujnik_podcisnienia_dodac}).
	\\
\end{easylist}

Na górnej warstwie pomiędzy złączami znajduje się dużo wolnej przestrzeni na której nie ma umiejscowionych żadnych komponentów. Jest to obszar rozbiegowy dla elastycznej tasiemki łączącej PCB z płytką zasilania (sekcja \ref{ss:power_board} i rysunek \ref{grafiki/glowica.png} c)).

\insertImgSetSize{grafiki/plytka_glowna_warstwy.png}
	{155}
	{Poszczególne warstwy płytki głównej sterownika, od lewej górna, dolna, warstwa zasilania ({\it power plane}), warstwa masy ({\it ground plane})}
	{mechsyspcb}

W następnych sekcjach tego rozdziału zostaną szczegółowo opisane poszczególne moduły elektroniczne płytki głównej. Podobnie jak w wypadku PCB, schematy ideowe tej płytki są o wiele obszerniejsze niż prezentowane wcześniej. Dlatego zostaną zaprezentowane tylko kawałki schematów niezbędne do opisu poszczególnych funkcjonalności. Kompletny schemat płytki znajduję się w załączniku \ref{dac zalacznik}. 

\insertImgSetSize{grafiki/top_level.pdf}
	{115}
	{Schemat ideowy najwyższego poziomu łączący ze sobą poszczególne moduły w projekcie płytki, w programie Altium Designer}
	{mechsyspcb}
	
Struktura schematu ideowego ma charakter obiektowy- podobnie jak w programowaniu wysokopoziomowym, zdeklarowane są klasy i tworzone są ich egzemplarze. Pozwala to na enkapsulacje modułów niskiego poziomu i powielanie ich w prosty sposób tylko na PCB, co przyspiesza projektowanie i zwiększa czytelność schematów.

\subsubsection{Zasilanie}

W układzie sterownika potrzebne jest wiele poziomów napięć ze względu na różnorodność zastosowanych układów scalonych. Aby je wszystkie uzyskać zostały zastosowane różne rozwiązania, zarówno bazujące na regulatorach impulsowych jak i liniowych. Wszystkie tory napięciowe są poprowadzone na jednej z warstw wewnętrznych PCB. \\

Do głowicy jest doprowadzone tylko jedno napięcie, z którego uzyskiwane są wszystkie pozostałe, niższe poziomy. To napięcie o amplitudzie 24 V samo w sobie jest także używane jeszcze przed redukcją do zasilania uzwojeń obu silników. ,,Logika'' urządzenia jest na poziomie 3.3 V i z tego względu większość elementów na płytce zasilana jest właśnie tym napięciem. Część układów natomiast wymaga do działania napięcia 5 V- W celu efektywnego uzyskania obu napięć został użyty obwód z rys. \ref{grafiki/power_5v_3v.pdf}.

\insertImgSetSize{grafiki/power_5v_3v.pdf}
	{52}
	{Część schematu zasilania prezentująca sposób uzyskania napięć 5 i 3.3V na płytce sterownika}
	{mechsyspcb}

W pierwszej kolejności napięcie główne trafia na przetwornicę A8498 (Allegro Microsystems) typu {\it Buck-converter}, która redukuje je do 5V na wyjściu. Układ jest w stanie dostarczyć prąd nawet do 3A, ale ze względu na niewielką przestrzeń dostępna na płytce elementy mocy przetwornicy (Cewka L1 i dioda D6) zostały dobrane tak, aby być w stanie dostarczyć tylko 400 mA prądu. Natężenie to jest w zupełności wystarczające do zasilania układów logicznych i analogowych sterownika, także tych pracujących przy napięciu 3.3 V- do wyjścia przetwornicy podłączony jest liniowy stabilizator napięcia LM3940 (Texas Instruments), który też pobiera prąd z przetwornicy.

\insertImgSetSize{grafiki/power_12_25.pdf}
	{60}
	{Źródła napięcia odniesienia dla obwodu analogowego sterownika}
	{mechsyspcb}
	
Urządzenia LM4040 i LM4041 (rys. \ref{grafiki/power_12_25.pdf}) to diody zenera pracujące jako precyzyjne układy napięcia odniesienia, których napięcie wyjściowe jest odporne na wpływ zmian temperatury otoczenia. Są używane do zasilania układów analogowych stosowanych w obwodzie sterowania silnikami. Więcej na ten temat jest zawarte w sekcji \ref{referencja}. Ostatnie napięcie to 12V, które jest także uzyskiwane z przetwornicy A8498, w nieco zmienionej konfiguracji (rys. \ref{grafiki/power_12v.pdf}).

\insertImgSetSize{grafiki/power_12v.pdf}
	{60}
	{Przetwornica typu Buck-converter dająca na wyjściu 12V, zasilająca sterowniki typu {\it High-Side}}
	{mechsyspcb}

Zmianie względem poprzedniego obwodu z tą przetwornicą uległy wartości rezystorów tworzących dzielnik napięcia między wyjściem przetwornicy a jej wejściem FB ({\it Feedback}), oraz połączenie do wejścia VBIAS, którego na tym schemacie już nie ma. Rezystory R9/ R11 na schemacie \ref{grafiki/power_12v.pdf} i R1/ R3 na schemacie \ref{grafiki/power_5v_3v.pdf} ustalają wartość napięcia wyjściowego głowicy. Zwierając VBIAS do wyjścia można zasilić układ logiczny w środku A8498. Wtedy wewnętrzny układ zasilający logikę może być wyłączony i chip wydziela mniej ciepła. Niestety można tak zrobić tylko do napięcia 5 V. \\

Napięcie wyjściowe przetwornicy jest potrzebne do zasilania bramek górnych tranzystorów tworzących mostki H w obwodach silników (sekcja \ref{dacref}). Amplituda tego napięcia może wynosić od 10 do 20 V, gdyż jest to wartość $ V_{GS} $ (Napięcie między bramką a źródłem) zastosowanych tranzystorów mocy przy której ich $ R_{DS(ON)} $ (Rezystancja przewodzenia między drenem a źródłem) jest bardzo niskie. Przyłożenie powyżej 20 V do bramki tranzystora może spowodować jego uszkodzenie. Dużą szkodą jest toz że napięcie doprowadzane do głowicy nie jest mniejsze bądź równe 20 V- w tym wypadku można by zrezygnować z tej przetwornicy i zasilać obwód sterujący silnikami bezpośrednio. Wartość tego napięcia wynika z doboru pozostałych elementów w maszynie. \\

Jeśli w przyszłych wersjach sterownika zajdzie potrzeba wygospodarowania dodatkowego miejsca na płytce głównej, to obwód z rys. \ref{power_12v.pdf} będzie można zastąpić stabilizatorem liniowym, dającym na wyjściu ok trochę poniżej 20 V. Taki stabilizator zajmowałby fizycznie mniej miejsca, ale w zamian wydzielała by się na nim większa moc, przez co generowałby więcej ciepła:

\begin{equation} \label{eq:zasilanie1}
	P = (V_{we} - V_{out})I
\end{equation}

Rys. \ref{grafiki/pcb_zasilanie.png} prezentuje model 3D części płytki na której widać oba omówione obwody zasilania. Aby maksymalnie wykorzystać dostępną przestrzeń na PCB, moduły zostały między sobą przedzielone obwodem sterowania silnikami.

\insertImgSetSize{grafiki/pcb_zasilanie.png}
	{36}
	{Rozmieszczenie modułów zasilania na płytce głównej, A: obwód \ref{grafiki/power_12v.pdf}, B: obwód \ref{grafiki/power_5v_3v.pdf}}
	{mechsyspcb}
	
\subsubsection{Jednostka centralna}

Głównym elementem sterującym peryferiami w urządzeniu jest 32 bitowy mikrokontroler STM32407VGT6 z rodziny Cortex® M4 (rodzina ARM) produkowany przez firmę ST Microelectronics. Kluczowymi parametrami decydującymi o doborze układu sterującego były:

\begin{easylist}
	& częstotliwość taktowania,
	& posiadanie jednostki zmiennoprzecinkowej ({\it FPU}- Floating Point Unit),
	& peryferia wewnętrzne pozwalające na sterowanie dwoma silnikami krokowymi (najwięcej końcówek mocy),
	& cena.
	\\
\end{easylist} 

Używając wspomnianego mikrokontrolera udało się zaimplementować wszystkie dotychczasowo działające funkcjonalności i na razie nie są znane żadne przeciwwskazania, które mogłyby by się pojawić w przyszłości. Nie było to jednak prostym zadaniem- urządzenie jest typowym procesorem ,,ogólnego użytku'' i nie posiada zaawansowanych peryferiów tak bardzo przydatnych przy sterowaniu silnikami, jak procesory ku temu dedykowane. 

\insertImgSetSize{grafiki/stm32f407vgt6.png}
	{45}
	{Mikrokontroler STM32F407VGT6 w obudowie TQFP100 o wymiarach 14 x 14 mm}
	{xiawu}
	
W dalszej części tego rozdziału zostały opisane w szczegółach sposoby w jaki poradzono sobie z ograniczeniami mikrokontrolera. 

\insertTab{|c|c|c|}
{%
\hline Mikrokontroler & STM32F407VGT6 & TMS320F28377S  \\
\hline Rodzina & ARM & C28x/ CLA    \\
\hline Architektura & RISC\footnote{{\it Reduced instruction set computing}} & RISC \\
\hline FPU & Single point & Single point \\
\hline \makecell{Maksymalna \\ częstotliwość \\ taktowania} & 168 Mhz & 200 Mhz \\
\hline MIPS \footnote{{\it Million Instructions Per Second}} & 210 & 400 \\
\hline \makecell{Pamięć \\ flash} & 1 MB & 1 MB \\
\hline \makecell{Pamięć \\ RAM} & 194 KB & 164 KB \\
\hline Ilość rdzeni &  1 &  2 \footnote{1 główny + {\it ,,control co-processor''} do obsługi procedur sterowania}  \\
\hline Cena \footnote{Przy zamówieniu 100 sztuk w hurtowni Digi-key} & \$ 9.34 & \$ 22.15 \\
\hline
}
{Porównanie mikrokontrolerów STM32F407VGT6 (ST Microelectronics) i TMS320F28377S  (Texas Instruments)}
{oprWlasne}
{tab:mcu_porownanie}

W tabeli \ref{tab:mcu_porownanie} przedstawiono porównanie zastosowanego układu z mikrokontrolerem dedykowanym pod sterowanie silnikami z najwyższej półki. 

\insertImgSetSize{grafiki/pcb_mcu.png}
	{65}
	{Widok 3D mikrokontrolera sterującego i otaczających go komponentów}
	{mechsyspcb}

Mimo że STM pod względem wymienionych parametrów zdecydowanie przegrywa, to ma jedną ważną zaletę, przez którą zadecydowano o jego wyborze- Jego współczynnik ceny do możliwości jakie posiada jest bardzo wysoki, możliwe że nawet jest jednym z najwyższych na rynku wśród mikrokontrolerów podobnej mocy. Ponadto jest łatwo dostępny, w większości przypadków ,,od ręki''. Jest to spowodowane tym że na jego bazie powstała płytka ewaluacyjna {\it STM32-DISCOVERY}, która cieszy się bardzo wysoką popularnością wśród hobbystów i profesjonalistów. \\

Układy STM32 są bardzo proste w programowaniu i debugowaniu dzięki dostępnym interfejsom {\it JTAG} i {\it SWD} (złącze CON3 na rys. \ref{grafiki/pcb_mcu.png}). JTAG jest interfejsem szybszym i bardziej stabilnym od SWD, ale wymaga większej ilości połączeń. W tym projekcie nie mógł zostać zastosowany gdyż część z pinów, które są potrzebne do jego pracy została użyta przy innych modułach. Źródłem zegarowym dla układu jest zewnętrzny mikro oscylator (Y1) o częstotliwości 16 Mhz. Używając wewnętrznych pętli PLL ({\it Phase Locked Loop}) uzyskiwana jest główna częstotliwość taktowania procesora 168 Mhz. Rysunek \ref{grafiki/clock_pll.png} przedstawia konfigurację wewnętrznych zegarów w mikrokontrolerze. Szybkie działanie mikrokontrolera jest parametrem kluczowym dla konstruowanego sterownika ze względu na to że silniki są maszynami wymagającymi obsługi w czasie rzeczywistym, a ponadto procesor musi jeszcze dodatkowo wykonywać setki obliczeń.

\insertImgSetSize{grafiki/clock_pll.png}
	{65}
	{Konfiguracja pętli PLL dla częstotliwości maksymalnej 168 Mhz}
	{oprWlasne}

W celu zapewnienia filtracji zakłóceń na liniach zasilania powodowanych głównie pracą silników, każde z wejść zasilających mikrokontrolera jest odprzęgnięte poprzez kondensator 100 nF (w obudowach 0402 tak jak większość komponentów pasywnych na płytce). \\

Dodatkowo na rysunku \ref{grafiki/pcb_mcu.png} widać złącza CON4 (z lewej) i CON7 (na dole), które łączą płytkę główną kolejno z enkoderem magnetycznym (\ref{ss:encmag}) i optycznym (\ref{ss:encopt}). Połączenie z tym drugim jest zrealizowane tasiemką tego samego typu co doprowadzane jest zasilanie do głównej PCB.








\clearpage


















\clearpage

















% Opis oprogramowania
\section{Oprogramowanie}
\label{s:oprogramowanie}

Program wykonywany w mikroprocesorze sterującym został napisany w języku c i w nielicznych częściach w assemblerze. Jądrem całej architektury jest system operacyjny czasu rzeczywistego FreeRTOS™. Wszystkie procedury wysokiego poziomu wykonywane w kodzie są umieszczone w wątkach o odpowiednich priorytetach, które porozumiewają się między sobą za pomocą kolejek, semaforów oraz obszarów pamięci współdzielonej. Zadania najniższego poziomu są wykonywane w procedurach obsługi przerwań, np. komutacja lub kolejkowanie ramek CAN. \\

Rysunek \ref{grafiki/Soft_architecture.pdf} przedstawia poglądowy schemat blokowy architektury systemu. Pomarańczowymi blokami zostały oznaczone procedury wykonywane na poziomie sprzętowym, lub programowym na warstwie znajdującej się niżej od wątków systemu FreeRTOS™. Kolor zielony oznacza wątki/ procesy w systemie FreeRTOS™. Strzałki z paskami symbolizują komunikaty przychodzące do kontrolera ze świata zewnętrznego, a bez nich komunikaty wewnątrz systemowe. Dokładniejszy opis poszczególnych bloków znajduje się w korespondujących z nimi sekcjach.

\insertImgSetSize{grafiki/peryferia.png}
	{110}
	{Widok wstępnie skonfigurowanych peryferiów w programie mxCube firmy ST Microelectronics}
	{oprWlasne}
	
Do wygenerowania kodu zawierającego wstępną konfigurację peryferiów służy program mxCube firmy ST Microelectronics.
	
\insertImgSetSize{grafiki/Soft_architecture.pdf}
	{220}
	{Schemat blokowy architektury oprogramowania sterownika}
	{oprWlasne}
	
\clearpage

Jego możliwości konfiguracyjne są bardzo podstawowe, ale pozwala on w graficzny sposób zobrazować używane moduły i wyprowadzenia (rys. \ref{grafiki/peryferia.png}). Weryfikację symulacyjną i analizę graficzną działania złożonych algorytmów przeprowadzono w środowisku MATLAB/ SIMULINK.

\subsection{Komunikacja}

Gotowe do obsługi programowej w urządzeniu są dwie magistrale szeregowe (patrz \ref{sss:hardware_komunikacja}). W tej sekcji znajduje się opis zastosowania każdej z nich na obecnym poziomie rozwoju.

\subsubsection{Magistrala CAN}
\label{ss:hwcan}

Magistrala pracuje w protokole CAN 2.0B (rozszerzona ramka: 29 bitowy identyfikator i 64 bitowe pole danych) z prędkością 1 Mbit/s. Używany protokół to TML CAN, który jest używany przed dotychczas używane w maszynach sterowniki firmy Technosoft. Jak zostało wcześniej wspomniane, projektowane urządzenie ma za zadanie zastąpić w maszynie dwa dotychczasowe, dlatego z punktu widzenia sieci są to dwa kontrolery (rys. \ref{grafiki/can_nodes.pdf}).

\insertImgSetSize{grafiki/can_nodes.pdf}
	{60}
	{Topologia sieci CAN. Sterownik przedstawia się jako dwa osobne urządzenia}
	{st}
	
W rzeczywistości kontroler jest tylko jeden. Mikrokontroler STM32F07VTG6 posiada dwie sprzętowe kolejki typu FIFO ({\it First In First Out}) do odbierania ramek CAN, które są współdzielone przez moduły CAN1 i CAN2. Urządzenie korzysta tylko z jednego peryferia, więc obie kolejki mogą być zastosowane przez nie. Filtry ramek przychodzących są skonfigurowane w taki sposób, aby przepuszczać tylko te ramki, które posiadają w swoim identyfikatorze numer osi Z lub R (odpowiednio 5 i 6). Po odebraniu wiadomości dla osi Z, ramka trafia do pierwszej kolejki FIFO, jeśli wiadomość jest dla osi R, to do drugiej. Taka architektura zapewnia bardzo dobre odseparowanie wirtualnych modułów odbiorczych sterowników już na poziomie sprzętowym. Bufory FIFO mają jednak ograniczoną pojemność-- mieszczą tylko 3 ramki. Z tego względu wiadomości do nich wpadające muszą być jak najszybciej odczytane i ulokowane w programowej kolejce o rozmiarze 32 ramek, w której mogą oczekiwać na interpretację przez dłuższy okres. \\

Dostęp do zmiennych (obszarów pamięci) w sterownikach Technosoft'u odbywa się poprzez zapis do i odczyt z rejestrów. Architektura konstruowanego kontrolera nie wspiera w natywny sposób takiego programowania, dlatego w celu poprawnej pracy potrzebna jest dodatkowa warstwa (rys. \ref{grafiki/tml_can_trans.pdf}).

\insertImgSetSize{grafiki/tml_can_trans.pdf}
	{35}
	{Sposób tłumaczenia ramek TML CAN w kontrolerze. Przykładowy zapis prędkości ruchu w osi Z}
	{oprWlasne}
	
W warstwie tłumaczenia ramek TML dokonywanych jest wiele różnych operacji. W zależności od komunikatu może to być zapis do zmiennej lub np. wywołanie funkcji. W przykładzie zaprezentowanym na rysunku \ref{grafiki/tml_can_trans.pdf} zewnętrzne urządzenia wysyła do kontrolera wiadomość TML, która po przeskalowaniu ustawia prędkość ruchu w osi Z.

\subsubsection{Magistrala RS232}
\label{ss:hwrs232}

Docelowo interfejs RS232 ma umożliwiać alternatywną metodę komunikacji \linebreak z urządzeniem w celu programowania i sterowania. Obecnie jednak jest to kanał typowo serwisowy. Korzysta z niego moduł logera (rys. \ref{grafiki/Soft_architecture.pdf}) aby wysyłać do aplikacji komputerowej informacje o zdarzeniach w HRF ({\it Human Readable Format}). Typowa ramka tego typu zawiera informacje o czasie wystąpienia zdarzenia, poziomie logowania (np. informacja o priorytecie normalnym lub krytycznym) oraz tekst dowolny (listing \ref{kody/rs232.txt}). 

\insertCode{kody/rs232.txt}
		   {}
		   {Przykładowe logi informujące o wywołaniu funkcji poprzez CAN}
		   {oprWlasne}
		   
Innym trybem logowania jest format CSV ({\it Comma Separated Values}), który pozwala na zbieranie danych w czasie rzeczywistym w celu ich dalszej analizy w programach MATLAB lub Excel. Obsługa modułu logera nie pochłania dużej ilości czasu procesora ze względu na to, że wątek ściągający kolejne logi z kolejki oraz kanał DMA, przez które są one wysyłane mają najniższe priorytety-- procedury wykonują się tylko wtedy, kiedy pozostałe procesy są uśpione.



\clearpage









\subsection{Realizacja zadań}

Wysunięcie się cylindra silnika VCM o odległość $ x $, zwrócenie wartości podciśnienia w głowicy czy obrót komponentu o 90\degree -- wszystko to jest przykładam wywołania funkcji, czy też realizacją zadania. Każda funkcja ma na wykonanie się określoną ilość czasu i jest zbudowana na zasadzie skryptu, który czeka na zajście konkretnych wydarzeń ({\it Events}), interpretuje je i wykonuje kolejne procedury w zależności od ich wyniku. W przypadku nie wykonania się którejkolwiek części skryptu w zadanym czasie, skrypt jest anulowany i do urządzenia wydającego komendę zostaje zwrócony kod błędu. System interpretuje obecnie 9 zdarzeń (listing \ref{kody/sys_triggers.h}).

\insertCode{kody/sys_triggers.h}
		   {C}
		   {Obiekt typu {\it enum} prezentujący obsługiwane w sterowniku wydarzenia}
		   {mechsyscode}

Zdarzenia mogą zachodzić sekwencyjnie lub równolegle. Skrypt może oczekiwać na zajście kilku skryptów (operatory {\it AND} i {\it OR}), gdzie każdy z nich także posiada ustalony czas na wykonanie się. \\

Rysunek \ref{grafiki/uml_homing.pdf} przedstawia przykładową procedurę (homing osi Z) w języku UML ({\it Unified Modeling Language}) wykonywaną przez system. W pierwszej kolejności zewnętrzne urządzenie (Kontroler systemowy automatu Pick and Place) wysyła poprzez magistralę CAN wywołanie funkcji Homing'u. Ramka trafia do kolejki odbiorczej CAN. Po interpretacji zostaje wybrany odpowiedni skrypt ({\it JobDoHoming}) i rozpoczyna się jego wykonywanie. \\

Pierwszym etapem działania skryptu jest zamknięcie zaworu podciśnienia. Następnie sprawdza na którym polu paska referencyjnego znajduje się część ruchoma głowicy (sekcja \ref{ss:encopt}). Jeśli jest to pole czarne (cylinder silnika VCM wsunięty), następuje wejście w pierwsze rozwidlenie w skrypcie (Black strip). Procedura homing'u służy do kalibracji pozycji. Zanim zostanie wykonana, absolutna pozycja cylindra nie jest znana. Po włączeniu zasilania przyjmuje się że pozycja w której aktualnie znajduje się cylinder to zero. Wiedząc tylko że cylinder jest w polu czarnym paska, następuje wysuw cylindra według określonej trajektorii (sekcja \ref{sss:trajgen}). Od razu po tym, funkcja oczekuje na zejście z czarnej strefy paska na białą (wiadomość {\it EXTI} na diagramie). Kiedy to nastąpi, zostaje ustalona absolutna pozycja cylindra. Od tej chwili można swobodnie sterować wysuwem głowicy. Na końcu cylinder wraca do ustalonej pozycji startowej.

\insertImgSetSize{grafiki/uml_homing.pdf}
	{130}
	{Diagram interakcji prezentujący wykonanie przez system procedury bazowania ({\it homing}) w osi Z}
	{oprWlasne}

Jeśli z jakiegoś powodu podczas uruchomienia procedury homing'u cylinder nie jest wsunięty (np. rozciągnięta lub uszkodzona sprężyna) i głowica znajduje się nad polem białym, najpierw następuje jej wsunięcie, aby homing zawsze następował \linebreak w ten sam sposób (zejście z pola czarnego na białe). Potwierdzenie wykonania zadania jest wysyłane do urządzenia wywołującego funkcję, w tym wypadku do kontrolera systemu. \\

Wszystkie zadania obecnie obsługiwane przez sterownik są wykonywane na prezentowanej na rys. \ref{grafiki/uml_homing.pdf} zasadzie (np. podniesienie komponentu, położenie go lub obrócenie). 






\subsection{Algorytmy enkodera magnetycznego}

W odróżnieniu do enkodera liniowego, dane odczytywane ze skonstruowanego modułu enkodera magnetycznego bazującego na układzie AS5048A (rys. \ref{grafiki/AS5048A_diagram.eps}) są obarczone błędami związanymi z nieliniowością układu i typem enkodera. Bez zastosowania żadnych algorytmów uśredniających odczytaną pozycję kątową, dokładność systemu była by marna, a sam układ napędowy działał by niestabilnie. \\

Aby umożliwić pracę enkodera z rozdzielczością do 0.05\degree{} w sterowniku zostały zaimplementowane następujące algorytmy:

\begin{easylist}
	& Uśrednianie bazujące na średniej arytmetycznej,
	& Filtr Kalmana (opracowany na podstawie \cite{kalman} i \cite{forbot}),
	& Filtr dolnoprzepustowy jednobiegunowy, rekursywny (opracowany na podstawie \cite{smith}).
	\\
\end{easylist} 

\subsubsection{Uśrednianie}

Pierwszy z wymienionych algorytmów wykonywany jest zawsze. Pozostałe dwa natomiast są uruchamiane naprzemiennie, w zależności od konfiguracji użytkownika. Uśrednianie działa na prostej zasadzie-- Wykonywanych jest kilka odczytów pozycji z enkodera (ilość odczytów zależy od konfiguracji, najczęściej cztery), z których następnie obliczana jest średnia arytmetyczna (wzór \ref{eq:alg1}). W kolejnym kroku odczyt jest normalizowany (wzór \ref{eq:alg2})-- mechanizm jest wymagany kiedy pozycja waha się między 359-tym a zerowym stopniem.

\begin{equation} \label{eq:alg1}
	pos = \frac{a_0 + a_1 + a_2 + ... + a_n}{n} 
\end{equation}

\begin{equation} \label{eq:alg2}
pos = \left\{
  \begin{array}{ll}
    pos + 360 & : pos < 0 \\
    pos - 360 & : pos > 360
  \end{array}
\right.
\end{equation}

\begin{easylist}
	& $ pos $ : absolutna pozycja enkodera [deg],
	& $ n $ : ilość próbek do uśrednienia.
	\\
\end{easylist} 

Zbyt duża ilość odczytów nie jest opłacalna, gdyż błąd odczytu przestaje znacząco maleć (wzór \ref{eq:alg3}).

\begin{equation} \label{eq:alg3}
	x = \frac{1}{\sqrt{n}}
\end{equation}

\begin{easylist}
	& $ x $ : szum enkodera (max. 0.06\degree{}),
	& $ n $ : ilość próbek do uśrednienia.
\end{easylist} 

\subsubsection{Filtr Kalmana}

Zastosowany algorytm jest metodą filtracji dynamicznej. Posługując się tym narzędziem, można wyznaczyć pomiarowo niedostępne zmienne jedynie na podstawie bieżących wartości wielkości pomiarowo dostępnych oraz znajomości modelu matematycznego łączącego ze sobą obydwie te grupy pomiarów. \\

Proces można przedstawić za pomocą dyskretnego modelu w przestrzeni stanu:

\begin{equation} \label{eq:alg4}
	x(t + 1) = Ax(t) + Bu(t) + v(t)
\end{equation}

\begin{equation} \label{eq:alg5}
	y(t) = Cx(t) + w(t) 
\end{equation}

Gdzie:
\begin{easylist}
	& $ x(t) $ : stan w chwili czasu $ t = 0, 1,... $,
	& $ y(t) $ : wyjście układu,
	& $ A $ : macierz stanu,
	& $ B $ : macierz wejścia,
	& $ C $ : macierz wyjścia,
	& $ v(t) $ : szum procesowy,
	& $ w(t) $ : szum pomiarowy,
\end{easylist} 

\insertImgSetSize{grafiki/przebieg_kf.eps}
	{110}
	{Przebieg czasowy prezentujący pozycję wału silnika krokowego: zadaną, mierzoną przefiltrowaną przez {\it KF} i  mierzoną nieprzefiltrowaną}
	{oprWlasne}

Zastosowany filtr jest dwuwymiarowy.Oznacza to że estymacji są poddawane dwa parametry: pozycja i prędkość. Model stanowy wygląda następująco:

\begin{equation} \label{eq:alg6}
	x = \begin{bmatrix}
       	pos \\ 
       	vel
     	\end{bmatrix}
\end{equation}

\begin{equation} \label{eq:alg7}
	A = \begin{bmatrix}
       	1 & -dt \\ 
       	0 & 1
     	\end{bmatrix}
\end{equation}

\begin{equation} \label{eq:alg8}
	B = \begin{bmatrix}
       	dt \\ 
       	0
     	\end{bmatrix}
\end{equation}

\begin{equation} \label{eq:alg9}
	C = \begin{bmatrix}
       	1 & 0 
     	\end{bmatrix}
\end{equation}

Implementacja dwuwymiarowa zapewnia dokładniejszą estymatę pozycji niż w przypadku kiedy pod uwagę brane było by tylko położenie. Jeszcze lepsze wyniki można by otrzymać stosując estymatę trójwymiarową (dodatkowo dodac parametr przyspieszenia). Niestety operacje na macierzach 3 x 3 pochłaniały za dużo zasobów procesora i z tego rozwiązania w sterowniku zrezygnowano. \\

Rys. \ref{grafiki/przebieg_kf.eps} prezentuje różne pozycje kątowe w funkcji czasu odczytane ze sterownika w trybie CSV (patrz sekcja \ref{ss:hwrs232}). Wstępnie widać że algorytm działa poprawnie, lecz jego zalety i wady można dokładniej zaobserwować na przebiegu z rys. \ref{grafiki/przebieg_kf_zoom.eps}.

\insertImgSetSize{grafiki/przebieg_kf_zoom.eps}
	{110}
	{Przebieg czasowy (powiększenie przy zmianie pozycji) prezentujący pozycję wału silnika krokowego: zadaną, mierzoną przefiltrowaną przez {\it KF} i  mierzoną nieprzefiltrowaną}
	{oprWlasne}
	
Po doświadczalnym dostosowaniu parametrów $ v $ i $ w $ algorytm tłumi zakłócenia bardzo dobrze. Charakterystyczną cechą filtru Kalmana jest jednak to że wraz ze wzrostem poziomu filtracji rośnie także opóźnienie sygnału. Parametr ten diametralnie wpływa na działanie regulatorów pozycji i prędkości (sekcja <dodac sekcje>). Pomimo tego że estymowana pozycja jest bardzo blisko zadanej, to czas potrzebny na jej osiągnięcie (w końcowym odcinku) jest zbyt długi. Z tego powodu został zaimplementowany dodatkowy, szybszy algorytm stosowany w czasowo krytycznych sytuacjach.

\clearpage

\subsubsection{Filtr dolnoprzepustowy}

Drugim algorytmem stosowanym w celu minimalizacji wpływu szumów towarzyszących odczytywanej pozycji kątowej z enkodera magnetycznego jest rekursywny (lub też {\it NOI}-- o nieskończonej odpowiedzi impulsowej), jednobiegunowy filtr dolnoprzepustowy. Zastosowane zostały jedynie dwa współczynniki rekursywne $ a_0 $ i $ b1 $:

\begin{equation} \label{eq:alg10}
	a_0 = 1 - x
\end{equation}

\begin{equation} \label{eq:alg11}
	b_1 = x
\end{equation}

Współczynnik $ x $ wskazuje jakie jest opadanie przebiegu mierzone między kolejnymi próbkami. Wartość ta jest obliczana na podstawie stałej czasowej filtru:

\begin{equation} \label{eq:alg12}
	x = \exp^{-2 \pi f_{CR}}
\end{equation}

$ f_{CR} $ jest częstotliwością odcięcia dobieraną w sterowniku doświadczalnie (obecnie 20 Hz). Kolejne próbki na wyjściu filtra obliczane są według wzoru \ref{eq:alg13}

\begin{equation} \label{eq:alg13}
	output(t) = input(t) a_0 + input(t - 1) b_1 
\end{equation}

\insertImgSetSize{grafiki/przebieg_lowpass_zoom.eps}
	{110}
	{Przebieg czasowy (powiększenie przy zmianie pozycji) prezentujący pozycję wału silnika krokowego: zadaną,   mierzoną nieprzefiltrowaną i mierzoną przefiltrowaną przez filtr dolnoprzepustowy trzeciego rzędu}
	{oprWlasne}
	
Rząd filtru jest ustalany przez użytkownika. Rys. \ref{grafiki/przebieg_lowpass_zoom.eps} prezentuje odczyt pozycji wału silnika krokowego. Algorytm ten działa szybciej od zaimplementowanego filtru Kalmana-- obrót o 90\degree{} w ok. 250 ms przy zastosowaniu KF i ok. 100 ms przy filtrze dolnoprzepustowym trzeciego rzędu. Ponadto, algorytm jest bardzo prosty i kod potrzebny do jego wykonania nie wymaga dużej ilości zasobów mikrokontrolera (listing \ref{kody/low_pass.c}).

\insertCode{kody/low_pass.c}
		   {C}
		   {Kod filtra dolnoprzepustowego wykonywany w pętli sterującej mikrokontrolera w interwale $ dt = 0.001 s $}
		   {mechsyscode}
		   
Na chwilę obecną uśrednianie wraz z  filtrem dolnoprzepustowym są algorytmami dające najlepsze wyniki działania enkodera.

















\subsection{Algorytmy sterowania wysokopoziomowego}





\clearpage

\subsection{Procedury sterowania niskopoziomowego}
\label{sss:lowlevel}

Cechą wspólną dla wszystkich rodzajów silników obsługiwanych przez sterownik jest zasilanie uzwojeń korzystając z  modulacji unipolarnej (sekcja \ref{sss:moduni}) i pomiar prądu tą samą metodą. 

\subsubsection{Modulacja szerokości impulsów}
\label{sss:pwm_soft}

Moduły PWM sterujące silnikami są skonfigurowanie w trybie symetrycznym ({\it Center aligned mode}). Oznacza to, że poszczególne kanały danego timer'a są względem siebie wyrównane do środka każdego impulsu (rys. \ref{grafiki/pwm_counter.eps}). Licznik w module inkrementuje i dekrementuje się co drugi okres. Fakt ten sprawia że maksymalna częstotliwość możliwa do uzyskania w tym trybie PWM jest dwa razy mniejsza od tej w trybie niesymetrycznym ({\it Edge aligned mode}). \\
	
Tryb symetryczny pozwala jednak na implementację modulacji unipolarnej \linebreak w bardzo prosty sposób. Na rys. (rys. \ref{grafiki/pwm_counter.eps}) widać przebiegi napięcia po obu stronach uzwojenia (A i B). W prezentowanym przykładzie wypełnienie impulsu przebiegów nie jest sobie równe, co oznacza, że element wykonawczy porusza się. 

\insertImgSetSize{grafiki/pwm_counter.eps}
	{100}
	{Potencjały na zaciskach uzwojenia wraz z licznikiem danego timer'a}
	{oprWlasne}

W momentach, w których potencjał A ma inną wartość niż B, prąd płynie przez uzwojenie poprzez naprzeciwnie otwarte tranzystory mostka H. Kiedy potencjały mają tą samą wartość, otwarte są te same tranzystory po obu stronach mostka (dwa górne lub dolne). Wtedy prąd przepływa przez uzwojenie, ale rozładowuje się. Aby przeliczyć zadany przez regulator prąd na wartość wypełnienia impulsu, w kontrolerze używane są następujące równania: 

\begin{equation} \label{eq:lowlevel1}
	K_{LR} = \frac{L_w F_{pwm}}{2 R_w}
\end{equation}

Równanie \ref{eq:lowlevel1} prezentuje wzór na stałą prądową uzwojenia. Poszczególne zmienne to:

\begin{easylist}
	& $ L_w $ : indukcyjność uzwojenia,
	& $ R_w $ : rezystancja uzwojenia,
	& $ F_pwm $ : częstotliwość przełączania kluczy,
	\\
\end{easylist} 

Dwójka we wzorze jest uproszczeniem, według którego przyjmuje się że w jednej stałej czasowej $ \tau $ obwodu $ RL $  prąd uzwojenia osiąga 50\% swojej maksymalnej wartości (\ref{eq:lowlevel2}), wynikającej z prawa Ohm'a dla obwodu w stanie ustalonym.

\begin{equation} \label{eq:lowlevel2}
	I_{max} = \frac{V_{CC}}{R_L + R_{meas}}
\end{equation}

Spadki napięć na tranzystorach są w równaniu \ref{eq:lowlevel2} pomijane. Stała $ K_{LR} $ jest następnie używana do obliczenia współczynnika PWM dla danego uzwojenia \ref{eq:lowlevel3}.

\begin{equation} \label{eq:lowlevel3}
	r_{pwm} = I_{reg} \cdot K_{LR} \cdot \frac{R_L + R_{meas}}{V_{CC}} \cdot \frac{T_{pwm}}{2}
\end{equation}

$ V_{CC} $ jest napięciem zasilania mostka, a $ T_{pwm} $ okresem PWM dla danego silnika. Na końcu obliczane są wartości wypełnień impulsów dla potencjałów po obu stronach uzwojenia \ref{eq:lowlevel4}.

\begin{equation} \label{eq:lowlevel4}
	PWM_{V1} = \frac{T_{pwm}}{2} + r_{pwm} 	\qquad 	PWM_{V2} = \frac{T_{pwm}}{2} - r_{pwm}
\end{equation}

Dla silnika krokowego obliczane są odpowiednio cztery wypełnienia impulsów zamiast dwóch.

\subsubsection{Pomiar prądu w uzwojeniu}

Rysunek \ref{grafiki/prady_H.eps} prezentuje uproszczony mostek H ze schematu \ref{grafiki/mostki_vcm.pdf}. Cyframi zostały oznaczone kolejne sekwencje komutacyjne, w których układ znajduje się przy sterowaniu unipolarnym (sekcja \ref{sss:pwm_soft}). Litery oznaczają sekwencję dla ruchu (np. cylindra w silniku VCM) w przeciwnych kierunkach. Strzałki wyznaczają drogę przepływu prądu i jego kierunek, natomiast blok w dolnej lewej gałęzi mostka symbolizuje rezystor pomiarowy (schemat \ref{grafiki/cursens_ls.pdf}). 

\insertImgSetSize{grafiki/prady_H.eps}
	{105}
	{Przepływ prądu w mostku H dla kolejnych komutacji (1, 2, 3, 4) przy ruchu elementu wykonawczego. a) Ruch w elementu wykonawczego w prawo, b) ruch w lewo}
	{oprWlasne}

Jak wynika z diagramu, prąd przepływa przez rezystor pomiarowy tylko w połowie możliwych konfiguracji. Fakt ten sprawia, że prąd nie może być mierzony \linebreak w dowolnym momencie, a wymaga synchronizacji z modułem timer'a, którego częścią są dane kanały PWM. Aby to zrobić, w kontrolerze posłużono się przerwaniami generowanymi w momencie, kiedy licznik danego timer'a zaczyna liczyć w dół lub w górę. Na rys. \ref{grafiki/pwm_counter.eps} są to momenty, w których wypełnienie impulsów jest w \linebreak 50\% swojej całkowitej szerokości. Te przerwania następnie uruchamiają przetwornik analogowo-cyfrowy (moduł skonfigurowany sprzętowo), który odczytuje wartość spadku napięcia na rezystorze. W przypadku silnika krokowego, jest to robione dwa razy (kolejno z obu rezystorów) ze względu na to że uzwojenia są dwa. \\

Rys. \ref{grafiki/stepper_prady_zoom.png} prezentuje oscylogram napięć odkładających się na rezystorach pomiarowych w obwodzie silnika krokowego. 

\insertImgSetSize{grafiki/stepper_prady_zoom.png}
	{140}
	{Oscylogramy mierzonych spadków napięć na rezystorach pomiarowych w mostku H silnika krokowego z rys. \ref{grafiki/stepper_prady.png} w przybliżeniu}
	{oprWlasne}
	
Można na nim zaobserwować poprawną pracę obwodu CSLS (schemat \ref{grafiki/cursens_ls.pdf}). \linebreak W co drugim okresie timer'a, spadki napięć na rezystorach pomiarowych wynoszą ok. 1.2 V co sugerowałoby, że w tym czasie przez uzwojenie nie płynie prąd (napięcie referencyjne). Powodem tego zjawiska jest jednak to, że tak jak pokazano na rys. \ref{grafiki/prady_H.eps} w połowie okresów prąd jest dla przetwornika ADC ,,niewidoczny''-- punkt pomiarowy jest dla niego w stanie wysokiej impedancji (nie ma różnicy napięć). \\

Mimo tego w oprogramowaniu mierzone są napięcia we wszystkich okresach, nawet tych niewidocznych. Implementacja mechanizmu synchronizującego przetwornik z okresami, zarówno jak jego działanie było by bardziej czasochłonne niż zastosowany obecnie algorytm. Mianowicie, w programie sprawdzane jest w którym okresie dla danej pary (wysoki--niski) amplituda jest większa i tylko ją stosuje się do obliczeń. Okres z napięciem referencyjnym jest odrzucany. \\

Prąd płynący w uzwojeniu jest liczony ze wzoru \ref{eq:lowlevel5}.

\begin{equation} \label{eq:lowlevel5}
	I_{meas} = \frac{V_{ref} \cdot ADC}{G \cdot RES_{ADC} \cdot R_{meas}}
\end{equation}

Gdzie:
\begin{easylist}
	& $ V_{ref} $ : napięcie źródła odniesienia,
	& $ ADC $ : wartość z przetwornika ADC,
	& $ RES_{ADC} $ : Rozdzielczość przetwornika (12 bitów),
\end{easylist} 




















\clearpage

% --------------------------------------------------------------------
% Bibliografia
% ustawienie bibliografii
\nocite{*}
\bibliographystyle{acm-i}
\bibliography{bib_database} 
\clearpage

\listoftables
\clearpage

\listoffigures
\clearpage

\lstlistoflistings




\end{document}

